\section{Laws for the corner state sum}\label{sec:knotlaws}

Throughout, $C$ is the function of Definition~\ref{def:C} on generic
polygons, and the named walls are those of Definition~\ref{def:walls}.

\begin{remark}[crossing records and knot presentation]\label{rem:crossing-record}
Three phrases of this section and of Section~\ref{sec:statesum} have the
following fixed meanings. (a)~The \emph{crossing record} of a diagram, here
always the positive lift of a carrier, is its named record in the sense of
Definition~\ref{def:gauss-record}: the crossing visits in their cyclic
traversal order, the pairing of the two visits of each crossing, the
over/under bit of each visit, and the crossing signs. For a positive lift
every crossing sign is $+1$ (Definition~\ref{def:positive-lift}), so the
record is the cyclic order, the pairing and the over/under bits; the words
\emph{order and crossing records} in the proof of Lemma~\ref{lem:flat-sides}
name the edgewise visit lists of the Gauss word
(Definition~\ref{def:gauss}), the part of this datum without the over/under
bits. (b)~The over/under bit is a chirotope datum: at a crossing of $E_i$
and $E_j$ the over strand of the positive lift is $E_i$ exactly when
$\det(\lt_i,\lt_j)>0$. On a chamber the crossing set is constant
(Proposition~\ref{prop:chambers}) and this determinant is continuous and
nonvanishing, the crossing being transverse, so its sign, and with it every
over/under bit, is constant on each chamber; the chamber-constancy and
silence proofs below print this argument where they use it. (c)~Two diagrams
\emph{present the same oriented knot} in the sense fixed by the statement of
Theorem~\ref{lit:gauss}: their compatible lifts in $S^3$ are joined by an
ambient isotopy preserving the traversal orientation. The only proof of this
document that uses this notion is the identification of the positive star
with a torus knot in Section~\ref{sec:anchors}, through that theorem. In
Literature input~\ref{lit:homfly} the phrase \emph{the oriented link presented
by $D$} is the source's notion, the oriented link type of the diagram, which
Reidemeister's theorem (registry L-1, Section~\ref{sec:registry}) carries
from the diagram.
\end{remark}

\begin{proposition}[chamber constancy]\label{prop:C-chamber}
The state sum $C$ of Definition~\ref{def:C} is constant on every chamber.
\status{proved (refereed: bench A, 2026-09-05T18:22:57Z; transcribed from CV prop:chamberinv, through Lemma~\ref{lem:carriers} and the named-record bridge (Lemma~\ref{rp:record-polynomial}, Theorem~\ref{lp:core}) (C020, C027); F-25-95; single convention, CV one-based tail = this document (F-25-104))}
\end{proposition}
\begin{proof}
First consider a continuous path of labelled generic polygons. By
Proposition~\ref{prop:chambers}, the crossing set and the order of its
visits along each directed edge are constant. Concatenating those edgewise
lists identifies the oriented cyclic traversal records, their crossing
pairings and interlacement graphs. It therefore gives a bijection of all
independent sets $S$, preserving $|S|$.

Fix one such support $S$. Include the original vertices as marked points
on the traversal circle and perform the same exchanges of outgoing
successors at its selected visits. By Lemma~\ref{lem:carriers}, with the
incoming-visit convention, this gives a fixed finite collection of oriented
carrier cycles. Corresponding original vertices and smoothing corners
have the same places on those cycles. An unselected crossing is retained
as a self-crossing of a cycle exactly when both its visits belong to that
cycle, so every carrier's retained crossing pairing and number $m_Q$ are
constant.

Each crossing point varies continuously: its two edge directions have
nonzero determinant, and solving the two-by-two intersection equations
gives continuous parameters in the relative interiors. Thus the corners
of each named carrier form a continuous ordered polygonal tuple.
Lemma~\ref{lem:carriers} puts it in the regular locus: its segments are
nonzero and no consecutive directions are antiparallel. Its signed
rotation is consequently constant by Lemma~\ref{lem:rot}(ii).
At an original corner the turn sign is the corresponding original turn;
at a smoothing corner it is the sign of the determinant of the inherited
ordered edge directions. These determinants are continuous and nonzero
along the path. All carrier corner signs, corner counts and uniformity
conditions therefore remain constant, as does the original left-turn
count $\ell(P)$.

At a retained self-crossing the two carrier directions are positive
multiples of the same ordered parent directions throughout. Their
nonzero determinant has constant sign. The physical over strand in the
positive lift therefore stays over, and the crossing sign stays positive.
Corresponding carrier lifts have a cyclic-order-, pairing-, sign- and
over/under-preserving bijection of crossing visits. The printed
scalar named-record Lemma~\ref{rp:record-polynomial}, followed by the
common substitution and Theorem~\ref{lp:core}, gives equal $H$ polynomials,
including for a carrier with no self-crossings.
Since $m_Q$ and $|\rot(Q)|$ are constant, the exponent
$d_Q=1-m_Q-|\rot(Q)|$ is constant. Hence the coefficient
$c(Q)=[a^{d_Q}z^0]H_Q^+$ is constant for every carrier, whether uniform
or mixed. Every term and the prefactor of Definition~\ref{def:C} are
now constant, proving constancy along the labelled path.

To pass to unlabelled chambers, first note that cyclically shifting the
vertex labels preserves $C$ directly from its definition. It merely
renames the same directed traversal, crossings and independent supports;
the corresponding carrier images and orientations, corner counts,
rotations, positive lifts and coefficients are unchanged. Openness of the
labelled generic locus gives each point a path-connected neighbourhood
inside it. The path argument makes $C$ locally constant there. Its
cyclic invariance makes it descend to a locally constant function on the
quotient by cyclic relabelling: the preimage of any value is an open
saturated subset, so its image is open in the quotient topology. A
locally constant function on a connected chamber is constant. This proves
the stated assertion without assuming a lift of an arbitrary quotient path.
\end{proof}

\begin{proposition}[silence]\label{prop:C-silent}
At a simple exterior-extension wall (E) or a simple pure cut (C),
\begin{equation}\label{ccf:silence}
C(P_+)=C(P_-).
\end{equation}
\status{proved (refereed: bench A, 2026-09-05T15:51:04Z; transcribed from exterior extension CV lem:silence; pure cut as printed (C020, C027))}
\end{proposition}
\begin{proof}
By Lemma~\ref{lem:wall-sides}(E),(C), the remote crossings and their
cyclic visit order agree on both sides and at the centre. Its proof also
gives the following geometric facts at the centre: vertices are distinct,
all edges and all original turns are nonzero, no vertex lies on a
nonincident closed edge, every actual crossing is transverse and interior,
and distinct crossings have distinct points. In (E) the sole vertex-line
incidence lies outside its closed base segment; in (C) its zero triple
contains no adjacent pair and hence gives no vertex-edge incidence.
Thus the same finite crossing records and independent supports are
identified through the entire small interval.

For each common support $S$, exchange the same successors to identify its
carriers. Original vertices and selected smoothing points remain distinct,
so the subsegments between successive marks have positive lengths at the
centre and on a smaller interval. At an original corner the two directed
segments are positive multiples of consecutive parent directions with
nonzero turn determinant. At a smoothing corner they are positive
multiples of the two transverse directions at its selected crossing,
again with nonzero determinant. Omitting unselected marks only joins
positive subsegments in the same direction. These observations verify
directly, including at the nongeneric parent centre, that every carrier
tuple is continuous, has nonzero segments, and has no antiparallel corner.
They do not apply a generic-polygon lemma outside its hypothesis.

It follows from Lemma~\ref{lem:rot}(ii) that each carrier's signed rotation
is the same on both sides. Every original and smoothing turn sign is
constant, and so are $\ell(P)$, uniformity and all corner counts. The
common successor cycles preserve which unselected crossing visits share
one carrier, hence all $m_Q$. At every retained crossing its direction
determinant is nonzero and continuous through zero; therefore its positive
over/under designation is also constant. Apply the scalar named-record equality of
Lemma~\ref{lc:presentations} and Theorem~\ref{lp:core} to corresponding
carrier diagrams on the two generic sides. This gives equal $H$ polynomials;
the equal $m_Q$ and $|\rot(Q)|$ give equality of their corner coefficients.
The independent supports, their uniformity, $|S|$ and the common prefactor
now identify the two sums term by term. There are finitely many supports
and carriers, so all interval shrinkings can be made simultaneously.

Only carrier geometry was used at the centre; the state sums in
\eqref{ccf:silence} are evaluated on the generic sides.
Proposition~\ref{prop:C-chamber} makes them independent of the chosen
representatives in the respective side chambers.
\end{proof}

\begin{theorem}[flat law]\label{thm:C-S3}
At a simple flat wall at $j$, with $n\geq4$,
\begin{equation}\label{ccf:flat-law}
C(P_{\rm right})-C(P_{\rm left})=C(P(0)\setminus j).
\end{equation}
\status{proved (refereed: bench A, 2026-09-05T15:51:04Z; transcribed from CV thm:main(A)(ii), lem:flatdata (C020, C027))}
\end{theorem}
\begin{proof}
We first rewrite the actual state sum using only the corner signs.
For a carrier $Q$ with $N_Q$ corners, define
\begin{equation}\label{ccf:selector}
W(Q)=
\begin{cases}
1,&\text{all its turns are right},\\
(-1)^{N_Q},&\text{all its turns are left},\\
0,&\text{its turns are mixed}.
\end{cases}
\end{equation}
If a support $S$ is uniform, the total number of left carrier corners is
$\ell(P)+|S|$: every original vertex belongs to exactly one carrier with
its original turn, and each selected crossing creates exactly one left
and one right smoothing corner, by Lemma~\ref{lem:carriers}. It follows
that the product of the selectors is $(-1)^{\ell(P)+|S|}$. If $S$ is
not uniform, at least one selector is zero. Consequently the definition
of $C$ gives, on every generic polygon,
\begin{equation}\label{ccf:selector-sum}
C(P)=\sum_{S\in\operatorname{Ind}(G_P)}
       \prod_{Q\text{ carrier of }S}W(Q)c(Q).
\end{equation}
This is an identity for the coefficients of the individual actual carrier
diagrams in Definition~\ref{def:C}; no factorization of one diagram's
polynomial into polynomials of other pieces has been used.

Put $D=P(0)\setminus j$. Lemma~\ref{lem:flat-sides} proves that $D$
is generic and identifies its crossings with those on both sides,
preserving cyclic order, pairing and positive over/under data. Hence it
identifies the three interlacement graphs and all independent supports,
including supports whose uniformity differs between the configurations.
For a fixed support, Corollary~\ref{cor:flat-carriers} gives corresponding
oriented carriers with the same retained crossing records and the same
signed rotations. Lemma~\ref{lc:presentations} and
Theorem~\ref{lp:core} therefore identify the $H$ polynomials of their
actual positive lifts. The numbers $m_Q$ and
$|\rot(Q)|$ agree, so the exponents $d_Q$ agree and each corresponding
coefficient $c(Q)$ is the same in all three configurations. This applies
also to mixed carriers and to crossing-free carriers.

Exactly one named carrier passes through the extra vertex $\mu_j$.
Call it $Q_*$ and call its corresponding deletion carrier $Q_*^D$.
All other corresponding selectors agree. The selectors of the
distinguished carrier satisfy
\begin{equation}\label{ccf:distinguished-selector}
W(Q_*^{\rm right})-W(Q_*^{\rm left})=W(Q_*^D)
\end{equation}
by Corollary~\ref{cor:flat-carriers}(iii). Explicitly, if the surviving
corners are mixed, all three entries are zero; if all are right, the
entries are $(1,0,1)$; if all $k$ surviving corners are left, the entries
are $(0,(-1)^{k+1},(-1)^k)$. In the last case the first minus the
second is $-(-1)^{k+1}=(-1)^k$. The deletion carrier has at least
three corners, so these alternatives do not involve an empty corner list.

For the fixed common support $S$, let $B_S$ be the common product of
all carrier coefficients, and let $U_S$ be the common product of all
selectors other than the distinguished one (with empty product $1$).
Its right-minus-left contribution to \eqref{ccf:selector-sum} is
\begin{equation}\label{ccf:term-difference}
B_SU_S\bigl(W(Q_*^{\rm right})-W(Q_*^{\rm left})\bigr)
=B_SU_SW(Q_*^D).
\end{equation}
The right side is exactly this support's contribution to the deletion
sum. No division by $B_S$, $U_S$ or a coefficient is performed; zeros
cause no exception. Summing over the finite common support set proves
\eqref{ccf:flat-law}. Proposition~\ref{prop:C-chamber} supplies the
well-defined side values. The only state sum on the right is that of
the generic deletion, not a value assigned to the flat centre.
\end{proof}

\begin{lemma}[products and the two-component row]\label{lem:homflyrows}
For oriented knots $K,J$, $H_{K\#J}=H_KH_J$ and
$H_{K\sqcup J}=\frac{a-a^{-1}}zH_KH_J$. For a two-component oriented link
diagram $D=D_1\cup D_2$ with linking number $\lk$,
$[z^{-1}]H_D=(a-a^{-1})a^{-2\lk}[z^0](H_{D_1}H_{D_2})$.
\status{proved (refereed: bench A, 2026-09-05T15:51:04Z; transcribed from CV lem:homflyrows via Theorems~\ref{mp:join}, \ref{mp:stack}, \ref{mp:lowest} (C028); its second sentence (the two-component row) consumes the knot clause of registry L-1 (F-25-48))}
\end{lemma}
\begin{proof}
Choose actual knot diagrams for $K,J$ with clean marked intervals. The
clean joining construction preceding Theorem~\ref{mp:join} gives a diagram
of their ordinary oriented connected sum. That theorem gives
$P_{K\#J}=P_KP_J$ on these actual representatives, and
Theorem~\ref{lp:core} identifies each $P$ with $H$. For split union, choose
representatives with disjoint page images. Theorem~\ref{mp:stack}, with
two one-component blocks and no mixed crossings, gives
$P_{K\sqcup J}=\delta P_KP_J$; the same identification gives the asserted
$H$ formula. Passing from these chosen representatives to the knot-class
notation uses the full global source premise explicitly retained in
Literature input~\ref{lit:homfly}; it is not a conclusion of the local
construction alone.

For the actual two-component diagram $D$, apply Theorem~\ref{mp:lowest}
with $c=2$. Its linking half-sum is computed from the original common
presentation, with the crossing signs of
Definition~\ref{def:positive-lift}, exactly the $\lk$ in the statement.
Its two intrinsic restrictions are the actual diagrams $D_1,D_2$.
Theorem~\ref{lp:core} identifies the result with
\[
[z^{-1}]H_D=(a-a^{-1})a^{-2\lk}[z^0]H_{D_1}\,[z^0]H_{D_2}.
\]
Both knot polynomials have only nonnegative even $z$ exponents by that
same theorem. Thus a product monomial has $z$ exponent zero precisely
when each factor does. It follows that the last two factors equal
$[z^0](H_{D_1}H_{D_2})$, which is the required expression. No division
by either knot polynomial or by a corner coefficient has been used.
\end{proof}

\begin{theorem}[vertex--edge law]\label{thm:C-S7}
At a simple vertex--edge wall at $(M;a)$, of bigon or sliding type, with
halves $\lambda_1,\lambda_2$ (Definition~\ref{def:deletion-halves}) and
contact sign $s=\chi_{a,a+1,M}(P_-)$,
\[
C(P_+)-C(P_-)=s\,C(\lambda_1)\,C(\lambda_2).
\]
\status{proved (refereed: bench A, 2026-09-05T15:51:04Z, 2026-09-06T06:11:53Z and 2026-09-06T07:18:25Z; transcribed from CV thm:s7universal(F) with its Section~6; RC d6\_vertexedge (C036))}
\end{theorem}
\begin{proof}
We use the actual carrier selector sum of Lemma~\ref{lem:C-X1}.
The geometric domain is exactly the simple SM vertex--edge wall in the
statement. Lemma~\ref{lem:wall-sides} supplies a clean contact disc,
stable old visits and the stated crossing changes. Lemma~\ref{lem:children}
supplies the actual generic ordered halves, each with at least three
vertices. No additional guarded-factor or parallel-direction hypothesis
is imposed.

In this proof write $\mathbf m=\mu_M$, $\mathbf a=\mu_a$ and
$\mathbf b=\mu_{a+1}$ for the positions along the germ, suppressing its parameter. Evaluations
at a specified side use that side's positions. The cyclic index $M$
in the theorem remains unchanged. A local diagram is always the actual
positive carrier diagram; its polynomial $P_A=H_A^+$ and writhe $m_A$
are those in the state sum. The blocks and their owners are those of
Definition~\ref{cb:blocks}; their products are proved in
Lemma~\ref{cb:products}. Equality of named scalar records uses
Lemma~\ref{rp:record-polynomial}, including crossing-free components.
Put $W(S)=\prod_L\operatorname{wt}(L)$ and let $\vartheta(u,v)$
denote the real principal angle from a nonzero vector $u$ to $v$.
All such local angles below are strictly between $-\pi$ and $\pi$.

\par\medskip\noindent\emph{Sliding: relocation and the exact selector difference.}\par

Put $r=\mathbf b-\mathbf a$, $u_{\rm in}=\mathbf m-\mu_{M-1}$ and
$u_{\rm out}=\mu_{M+1}-\mathbf m$ at the wall.  The two neighbours are on opposite
sides, while both tangents follow the oriented traversal through the
contact.  Therefore
\begin{equation}\label{s7c:sliding-signs}
 \sgn\det(r,u_{\rm in})=\sgn\det(r,u_{\rm out})\ne0.
\end{equation}
By continuity from the relative-interior contact, each incident line
meets the relative interior of the finite remote segment on a sufficiently
small event interval.  Exactly one incident segment reaches that line
intersection on each side; the two sides exchange which one.  Denote
the crossing on the chosen side $P_-$ by $x_-$ and its relocated mate
on $P_+$ by $x_+$.  Removing that label within the contact collar
leaves the same cyclic word.  Equation~\eqref{s7c:sliding-signs} also
fixes which visit is the remote overpass.  Thus the named map
$\varphi(x_-)=x_+$, identity on old labels, preserves cyclic order,
pairing, over/under, signs and the interlacement graph.  It transports
supports, undominated blocks, owners, actual carrier diagrams and writhes.

For a support avoiding this crossing, all successor cycles pass regularly
through the wall and all corner determinants stay nonzero.  Their
rotation integers and selectors are fixed, and so are their coefficient
reads.  This sector cancels term by term.

For a support containing $x_-$, the two visits cut the traversal into
the two half intervals.  Every other selected chord is internal to one
interval, by independence.  Conversely any two half supports together
with $x_-$ form a support.  This gives a bijection with explicit inverse
\begin{equation}\label{s7c:sliding-bijection}
 \{S\in\Ind(G_{P_-}):x_-\in S\}\longleftrightarrow
 \Ind(G_{\lambda_1})\times\Ind(G_{\lambda_2}),
 \quad S\longmapsto(S_1,S_2).
\end{equation}
Smoothing $x_-$ closes the two intervals separately.  All subsequent
selected smoothings are internal, and labels interlacing $x_-$ are
dominated.  Consequently
\begin{equation}\label{s7c:sliding-residual}
 G_{P_-}[U(S)]
 =G_{\lambda_1}[U(S_1)]\sqcup G_{\lambda_2}[U(S_2)].
\end{equation}
The named maps identify every residual diagram and every owner.  The
only changed direction list, on the carrier containing the short edge,
is one of
\begin{equation}\label{s7c:short-direction-lists}
 (r,u_{\rm in},u_{\rm out})\longrightarrow(r,u_{\rm out}),
 \qquad
 (u_{\rm in},u_{\rm out},r)\longrightarrow(u_{\rm in},r).
\end{equation}
The two incident tangents lie in one open angular half-plane bounded by
$r$.  Thus their principal turns satisfy the real equalities
\begin{align}
 \vartheta(r,u_{\rm in})+\vartheta(u_{\rm in},u_{\rm out})
       &=\vartheta(r,u_{\rm out}),\label{s7c:turn-short-a}\\
 \vartheta(u_{\rm in},u_{\rm out})+\vartheta(u_{\rm out},r)
       &=\vartheta(u_{\rm in},r).\label{s7c:turn-short-b}
\end{align}
There is no hidden multiple of $2\pi$.  The chamber and half carriers
have equal signed, and hence absolute, rotations.  Writing $C_\pm(S)$
for the complete coefficient product on the two sides, and $C_i(S_i)$
for its half versions, we get
\begin{equation}\label{s7c:sliding-coefficients}
 C_-(S)=C_+(\varphi S)=C_1(S_1)C_2(S_2).
\end{equation}

The half contact signs
$\eta_1=\sgn\det(r,u_{\rm out})$ and
$\eta_2=\sgn\det(u_{\rm in},r)$ are opposite.  In $P_-$ the
short-edge refinement is on the half with contact sign $s$, and in $P_+$
on the one with sign $-s$.  Let $h_\eta$ be that half-carrier weight,
$c_\eta$ its refined weight, and
$\tau=\sgn\det(u_{\rm in},u_{\rm out})$.  The refinement adds
the turn $\tau$, so the full selector table is
\begin{equation}\label{s7c:short-selector}
 c_\eta=\begin{cases}
 -\eta h_\eta,&\tau=\eta,\\
 0,&\tau=-\eta.
 \end{cases}
\end{equation}
Indeed a mixed original carrier stays mixed; if $\tau$ disagrees with
its contact sign it has both signs; if all signs are right the extra
right corner preserves weight $1$; and if all are left the extra left
corner reverses $(-1)^{\#\text{corners}}$.  These are all cases.
Let $Q_{\rm sp}$ be the product of the other carrier weights, without
assuming it nonzero.  For the row in \eqref{s7c:sliding-bijection},
\begin{equation}\label{s7c:sliding-selector-factors}
 W_1W_2=Q_{\rm sp}h_s h_{-s},\quad
 W_-=Q_{\rm sp}c_s h_{-s},\quad
 W_+=Q_{\rm sp}h_s c_{-s}.
\end{equation}
If $\tau=s$, then $c_s=-s h_s$ and $c_{-s}=0$; if $\tau=-s$,
then $c_s=0$ and $c_{-s}=s h_{-s}$.  Either way
\begin{equation}\label{s7c:sliding-selector-difference}
 W_+-W_-=sW_1W_2.
\end{equation}
Multiply by \eqref{s7c:sliding-coefficients} and sum over the finite
bijection.  Distributivity proves the stated product law for sliding,
with no spectator cancellation.

\par\medskip\noindent\emph{Bigon words and the two-newborn sector.}\par

Let $P_0$ be the side where $M$ and its two neighbours lie on the same
side of the remote line, and $P_2$ the other side, with two newborn
crossings $x,y$.  Set
\begin{equation}\label{s7c:bigon-signs}
 s_0=\sgn\det(r,\mathbf m-\mathbf a)\big|_{P_0},\qquad
 \delta_{\rm dir}=\begin{cases}+1&P_- =P_0,\\-1&P_-=P_2.
 \end{cases}
 \qquad s=\delta_{\rm dir}s_0.
\end{equation}
Write $\epsilon=1$ when the newborn chords interlace and $0$ otherwise.
Starting at the first visit of $x$, the two complete local words are
\begin{equation}\label{s7c:bigon-words-eq}
 \begin{array}{c|l}
 \epsilon=0&x_0y_0\,A\,y_1x_1\,B\\
 \epsilon=1&x_0y_0\,A\,x_1y_1\,B\\
 P_0&A\,B.
 \end{array}
\end{equation}
Each old label is internal to $A$, internal to $B$, or has one visit in
each.  The last set $N$ is the common old neighbourhood of $x,y$.
An old support $T$ is called eligible when $T\cap N=\varnothing$.
If $T$ meets $N$, both newborns are dominated and neither extension
by a newborn is a support.  Deleting their four visits identifies the
complete smoothing successor, owners, corners, rotations, carrier
diagrams and coefficient reads.  The entire ineligible returned sector
is therefore zero.

For eligible $T$, restriction to the intervals gives the bijection
\begin{equation}\label{s7c:eligible-bijection}
 T\longleftrightarrow(T_1,T_2)\in
 \Ind(G_{\lambda_1})\times\Ind(G_{\lambda_2}).
\end{equation}
Internal labels in the two intervals have disjoint induced graphs;
the halves close just these intervals.  Conversely the union of two
half supports is independent and misses $N$, proving the inverse.

If $\epsilon=0$, $T\cup\{x,y\}$ is a support.  Smoothing its four
newborn visits makes one local three-corner contact triangle and otherwise
exactly the successors of $T_1,T_2$.  Mixed old labels are dominated;
each unmixed residual label, owner and decorated diagram matches its half.
For $s_0=+1$ the triangle has three right turns; for $s_0=-1$ it has
three left turns.  Its complete data are
\begin{equation}\label{s7c:triangle-data}
 \rot=-s_0,\quad R=1,\quad P=1,\quad w=0,\quad
 [a^0z^0]1=1,\quad\operatorname{wt}=s_0.
\end{equation}
Every two-newborn support arises this way.  Put
$J=sC(\lambda_1)C(\lambda_2)$.  Finite distributivity and
\eqref{s7c:bigon-signs} show that the directed two-newborn sector is
$J$ when $\epsilon=0$.  When $\epsilon=1$ it is empty.  Hence, in
both branches,
\begin{equation}\label{s7c:sector-split}
 B=(1-\epsilon)J.
\end{equation}
It remains to show that the directed sum $R_{\rm ret}$ of all other
noncancelling rows equals $\epsilon J$.

\par\medskip\noindent\emph{Contact owners and one coambient smoothing diagram.}\par

Fix an eligible newborn-free support $T$.  Smooth its selected chords
inside $A$.  Choosing an innermost chord closes its inner subarc and
splices it out of the remaining boundary strand.  Induction leaves one
open boundary successor and internal cycles.  The same holds in $B$.
The two open boundary successors and contact stretches form one full
contact carrier $L^*$; closing them at the cut gives ordered half contact
carriers $L_1,L_2$.  All internal cycles are spectators.

For $i=1,2$, put $G_i=G_{\lambda_i}$ and
$U_i=V(G_i)\setminus(T_i\cup N_{G_i}(T_i))$.  Define the
owner-restricted sets
\begin{equation}\label{s7c:owner-restriction}
 O_i=\{c\in U_i:\text{the half owner of }c\text{ is }L_i\}.
\end{equation}
For an internal label, the successor correspondence gives full owner
$L^*$ if and only if its half owner is $L_i$.  Every surviving mixed
old label belongs to
$M_T=N\setminus N_{G_{P_2}}(T)$ and has owner $L^*$.
Thus the full high contact residual set is precisely
\begin{equation}\label{s7c:full-contact-set}
 O_1\sqcup O_2\sqcup M_T\sqcup\{x,y\}.
\end{equation}
The labels $U_i\setminus O_i$ are on spectator carriers and are never
inserted into a contact polynomial.

Let $L_L,D_L$ be the low full contact carrier and its actual positive carrier
diagram, and $L_H,D_H$ its high counterpart.  For this newborn-free support every selected crossing persists away
from the contact disc. Its successor carrier through the contact follows
the two original incident segments through $M$ on both sides. The edge
pieces have positive lengths at the centre, and all its corner turns
have nonzero determinants there. Thus this carrier is a continuous
regular family, including at the wall. Lemma~\ref{lem:rot}(ii) makes
its rotation locally constant. Consequently
\begin{equation}\label{s7c:full-rotation}
 \rot(L_L)=\rot(L_H)=\rot(L^*).
\end{equation}
The positive crossing $q=x$ of $D_H$ has oriented smoothing equal to
one actual ordered two-component diagram $D_A$.  Component 1 contains
the $A$-boundary successor and corresponds to $L_1$; component 2
contains the $B$-boundary successor and corresponds to $L_2$.
Retain every mixed crossing of this one ambient diagram.

If $\epsilon=1$ define $D_A^{\rm post}=D_A$.  If $\epsilon=0$,
component 1 contains the permitted positive self R-I curl $y$; delete
exactly that curl, retaining the two ordered tags and every mixed crossing,
to define $D_A^{\rm post}$.  Set
\begin{equation}\label{s7c:component-data}
 D_i=\operatorname{component}_i(D_A^{\rm post}),\quad
 Q_i=P_{D_i},\quad w_i=w(D_i),\qquad i=1,2.
\end{equation}
The successor words identify these as the actual carrier diagrams of
the ordered half contact carriers; their scalar identification is
Lemma~\ref{rp:record-polynomial}. The R-I deletion is in the supplied
clean curl disc. The local LM R-I equality applies to the entire
two-component diagram and to its component-1 restriction; component 2
is unchanged. It retains both component tags and every mixed crossing,
so the coambient mixed signed sum and hence linking number are unchanged.  Calling the latter $\ell$,
the full self/mixed crossing partition is
\begin{equation}\label{s7c:crossing-partition}
 \begin{array}{c|l}
 \text{component 1 self crossings}&O_1\cup(\{y\}\text{ if }\epsilon=0)\\
 \text{component 2 self crossings}&O_2\\
 \text{mixed crossings}&M_T\cup(\{y\}\text{ if }\epsilon=1).
 \end{array}
\end{equation}
In particular
\begin{align}
 P_{\operatorname{component}_i(D_A)}&=Q_i,\label{s7c:component-polynomial}\\
 \operatorname{lk}(D_A)&=\operatorname{lk}(D_A^{\rm post})=\ell,
 \label{s7c:ambient-linking}\\
 w(\operatorname{component}_1(D_A))&=w_1+1-\epsilon,
 \qquad w(\operatorname{component}_2(D_A))=w_2.
 \label{s7c:pre-curl-writhe}
\end{align}
Here $\operatorname{lk}(D_A)$ means the linking of its named ordered
components, not linking assigned to two separately embedded restrictions.

\par\medskip\noindent\emph{Turn and selector bookkeeping.}\par

Send $r$ to the positive horizontal ray.  If necessary reflect the
picture to place both neighbour vectors above it, with arguments
$\alpha,\beta\in(0,\pi)$.  Direct line intersection gives
\begin{equation}\label{s7c:angle-orders}
 \epsilon=0\ \Longleftrightarrow\ \alpha<\beta,
 \qquad
 \epsilon=1\ \Longleftrightarrow\ \beta<\alpha.
\end{equation}
The half contact turns are $\beta$ and $\pi-\alpha$; the full
contact turn is $\beta-\alpha-\pi$ in the noninterlacing case and
$\beta-\alpha+\pi$ in the interlacing case.  Both half signs are
$s_0$; the full sign is $-s_0$ or $s_0$, respectively.  Reflection
reverses all signs together and transports these statements back.
All noncontact turns partition between the halves.  If $Q_{\rm sp}$
is the common product of internal spectator selectors, the actual factors
are
\begin{equation}\label{s7c:selector-owners}
 W_T=Q_{\rm sp}\operatorname{wt}(L^*),\qquad
 W_1W_2=Q_{\rm sp}\operatorname{wt}(L_1)\operatorname{wt}(L_2).
\end{equation}
No factor is cancelled in these identities.

For interlacing the contact signs agree.  If any noncontact sign disagrees,
both the full weight and the relevant half product are zero.  Otherwise
all three carriers are uniform with the same sign.  For right turns each
weight is $1$ and $s_0=-1$; for left turns the full corner count is the
sum of the half counts minus one, giving a relative minus sign and
$s_0=+1$.  Thus, including all zero cases,
\begin{equation}\label{s7c:interlacing-selector}
 \operatorname{wt}(L^*)=-s_0\operatorname{wt}(L_1)
                       \operatorname{wt}(L_2),\qquad
 W_T=-s_0 W_1W_2.
\end{equation}
If $W_T\ne0$, \eqref{s7c:selector-owners} also shows that the spectator
product and all three contact weights are nonzero, so the just-described
uniform same-sign case applies.  Multiplication by a possibly zero
$Q_{\rm sp}$ proves the second identity without division.

For noninterlacing, a uniform full carrier gives each half one opposite
contact turn, while two uniform halves give the full carrier one opposite
turn.  Thus at least one of their selectors is zero:
\begin{equation}\label{s7c:noninterlacing-selector}
 W_T W_1W_2=0.
\end{equation}
Every carrier used here has at least three corners.  Two polygon corners
would force overlapping straight segments, a polygon and smoothing corner
would force vertex--edge incidence, and two smoothing corners would
contradict transversality of the crossing straight branches.  A one-corner
closed straight path is likewise impossible.  This also applies to the
daughter carriers used below.

Subtract the explicit full contact turn from the sum of the two half
turns, and add the unchanged noncontact turns.  Dividing by $2\pi$ gives
the signed identity
\begin{equation}\label{s7c:rotation-ledger}
 \rot(L_1)+\rot(L_2)-\rot(L^*)
 =\begin{cases}0&\epsilon=1,\\s_0&\epsilon=0.
 \end{cases}
\end{equation}
The second line contains a full-turn correction; unsigned angles would
not give this formula.

\par\medskip\noindent\emph{Universal skein extraction and the interlacing branch.}\par

For the fixed eligible row bind all polynomial and exponent objects:
\begin{align}
 F_L&=P_{D_L},&F_H&=P_{D_H},&F_A&=P_{D_A},\label{s7c:polynomial-binders}\\
 f_L&=[z^0]F_L,&f_H&=[z^0]F_H,&f_i&=[z^0]Q_i,\label{s7c:row-binders}\\
 w_L&=w(D_L),&w_H&=w(D_H),&R_i&=|\rot(L_i)|,\label{s7c:writhe-binders}\\
 R_L&=|\rot(L_L)|,&R_H&=|\rot(L_H)|,&&\label{s7c:rotation-binders}\\
 k_L&=1-w_L-R_L,&k_H&=1-w_H-R_H,&k_i&=1-w_i-R_i,\label{s7c:slot-binders}\\
 \Omega_L&=[a^{k_L}z^0]F_L,&\Omega_H&=[a^{k_H}z^0]F_H,
 &\omega_i&=[a^{k_i}]f_i,\label{s7c:coefficient-binders}\\
 K&=k_1+k_2.&&&&\label{s7c:product-slot}
\end{align}
These coefficient maps are total even when any Laurent polynomial is zero.
The two incident directions in this bigon branch have opposite
determinant signs against the remote direction. Their positive
decorations therefore put the remote strand over at exactly one
newborn crossing. Switching the positive contact crossing $q$ makes
the same strand over at both crossings, with opposite crossing signs.
The isolated contact disc contains no other strand; this is the actual
ordinary R-II template, not merely an opposite-sign pair. Delete it
using local LM invariance. The remaining diagram has the same complete
decorated record as $D_L$, including every old retained crossing;
Lemma~\ref{rp:record-polynomial} gives its value $F_L$. The actual
oriented smoothing is the tagged $D_A$. Thus the skein relation at $q$
first gives $aF_H-a^{-1}F_L=zF_A$, then
\begin{equation}\label{s7c:universal-skein}
 F_H=a^{-2}F_L+a^{-1}zF_A.
\end{equation}
The high writhe is $w_H=w_L+2$ and
\eqref{s7c:full-rotation} gives $R_H=R_L$, hence
$k_H=k_L-2$.  Extracting $[a^{k_L-2}z^0]$ in
\eqref{s7c:universal-skein} gives
$\Omega_H=\Omega_L+[a^{k_L-1}z^{-1}]F_A$.
Therefore
\begin{equation}\label{s7c:universal-extraction}
 \Omega_H-\Omega_L=[a^{k_L-1}z^{-1}]F_A.
\end{equation}
No selector or degree hypothesis is needed for this equality.

Assume first $\epsilon=1$.  The crossing partition
\eqref{s7c:crossing-partition} puts $y$ and all retained old mixed
crossings between the two components of $D_A$; none belongs to $Q_i$
or $w_i$.  Lemma~\ref{lem:homflyrows} and the writhe count yield
\begin{equation}\label{s7c:interlacing-writhe-row}
 [z^{-1}]F_A=a^{-2\ell}(a-a^{-1})f_1f_2,
 \qquad w_L=w_1+w_2+2\ell-1.
\end{equation}
The second equality follows also by counting $w_H$: the self crossings
contribute $w_1+w_2$, the mixed crossings of $D_A$ contribute $2\ell$,
and the smoothed crossing $q$ contributes one; subtract the two newborn
crossings to pass from high to low.

Let $U_{\rm sp}$ be the product of all spectator coefficient reads for
this newborn-free row.  Before using absolute-rotation or degree identities
tied to this row, test $W_T$ and $U_{\rm sp}$.  If $W_T=0$, the two
full terms vanish and \eqref{s7c:interlacing-selector} makes the half
selector product zero.  If $U_{\rm sp}=0$, both full terms and the
half product contain that same spectator zero.  Either case is finished.
Assume now $W_TU_{\rm sp}\ne0$.  The uniform same-sign implication
proved in the turn calculation above and
\eqref{s7c:rotation-ledger} imply
\begin{equation}\label{s7c:interlacing-absolute}
 R_L=R_1+R_2.
\end{equation}
Substitution of \eqref{s7c:interlacing-writhe-row} into the definition
of $k_L$, followed by \eqref{s7c:interlacing-absolute}, gives
\begin{align}
 k_L&=1-(w_1+w_2+2\ell-1)-(R_1+R_2)\notag\\
    &=(1-w_1-R_1)+(1-w_2-R_2)-2\ell
      =K-2\ell.\label{s7c:interlacing-slot}
\end{align}
In \eqref{s7c:universal-extraction}, the $a^{1-2\ell}$ term of
the mixed row requests degree $k_L+2\ell-2=K-2$ in $f_1f_2$;
its $-a^{-1-2\ell}$ term requests degree $k_L+2\ell=K$.
Consequently
\begin{equation}\label{s7c:interlacing-extraction}
 \Omega_H-\Omega_L=[a^{K-2}](f_1f_2)-[a^K](f_1f_2).
\end{equation}
If either $f_i=0$, both reads and $\omega_1\omega_2$ are zero;
no minimum degree is assigned.  Otherwise both half contact carriers
are uniform, and each is a subpolygon of a decomposition of its generic
half $\lambda_i$ (Lemma~\ref{lem:children}(ii), Lemma~\ref{lem:carriers}(i);
the hypothesis discharge below), the domain of Theorem~\ref{thm:floor},
which gives $\mindeg_a f_i\ge k_i$.
In the Laurent integral domain the product has degree at least
$k_1+k_2=K$.  Its degree $K-2$ coefficient is zero.  In the degree
$K$ convolution any nonzero pair has exponents at least $k_1,k_2$,
so the only possible pair is exactly $(k_1,k_2)$.  Thus
\begin{equation}\label{s7c:interlacing-coefficient-result}
 \Omega_H-\Omega_L=-\omega_1\omega_2.
\end{equation}

Both one-newborn sectors vanish by their selectors.  To check this
directly, take $\mathbf m=(0,0)$, $r=(1,0)$ on the line $y=c>0$, and
neighbour vectors $(u,h),(v,k)$ with $h,k>c$.  Interlacing means
$hv-uk>0$.  The corner determinants are
\begin{equation}\label{s7c:one-newborn-turns}
 \det((-u,-h),(v,k))=hv-uk>0,\quad
 \det(r,(-u,-h))=-h<0,\quad
 \det((v,k),r)=-k<0.
\end{equation}
The first turn is at $M$; the second or third is the smoothing turn
when $x$ or $y$ is selected.  Eligibility leaves the intervening collar
free of selected visits, so the opposite turns lie on the same daughter
carrier and make it mixed.  Reflection reverses all three signs and
proves the other orientation as well.

Let $\operatorname{Term}_H(T),\operatorname{Term}_L(T)$ denote complete
newborn-free support contributions, before multiplying by
$\delta_{\rm dir}$.  In the nonzero branch
\eqref{s7c:interlacing-coefficient-result} gives
\begin{equation}\label{s7c:interlacing-term}
 \operatorname{Term}_H(T)-\operatorname{Term}_L(T)
 =-W_TU_{\rm sp}\omega_1\omega_2.
\end{equation}
The same equation holds in the already discharged zero-factor branches,
with both sides zero.  Substitute $W_T=-s_0W_1W_2$, multiply by
$\delta_{\rm dir}$, and use $s=\delta_{\rm dir}s_0$.
The directed returned share is
$sW_1W_2U_{\rm sp}\omega_1\omega_2$.  The owner-restricted
decomposition puts each internal spectator into this half product
exactly once.  Summing over \eqref{s7c:eligible-bijection} gives
\begin{equation}\label{s7c:interlacing-return}
 R_{\rm ret}=sC(\lambda_1)C(\lambda_2)=J.
\end{equation}
This proves the required returned-sector identity for $\epsilon=1$.

\par\medskip\noindent\emph{Noninterlacing: full-twist, singleton and separate-block rows.}\par

Now $\epsilon=0$.  The two-newborn sector already equals $J$ by
\eqref{s7c:sector-split}; we show every remaining returned row is zero.

\par\noindent\emph{The same-block alternative.}\par
In $D_A$ the remaining newborn $y$ is a positive self curl on
component 1.  Its deletion gives the exact $Q_1,Q_2$ of
\eqref{s7c:component-data}, not polynomials with spectator factors
inserted.  Counting self and mixed crossings gives
\begin{equation}\label{s7c:noninterlacing-writhe}
 w_L=w_1+w_2+2\ell.
\end{equation}
Indeed $w_H$ includes the additional $q$ and the self curl $y$, for
a total of $w_1+w_2+2\ell+2$.
If the full selector or an actual spectator multiplier is zero, both
newborn-free terms vanish and the row stops before absolute-rotation
identities.  Otherwise the full carrier is uniform with sign
$\tau=-s_0$.  Every half inherits its old turns and has exactly one
new contact turn of sign $s_0$, hence exactly one dissent.  The
at-least-three-corner observation ensures inherited turns exist.
Lemma~\ref{lem:uniformrot}, reversing all orientations when
$\tau=-1$, puts the full and both half rotations on the $\tau$ ray.
Multiplying \eqref{s7c:rotation-ledger} by $\tau$ yields
\begin{equation}\label{s7c:noninterlacing-rotation}
 R_1+R_2-R_L=\tau s_0=-1.
\end{equation}
Let $\Delta_R=R_1+R_2-R_L$.  The slot definitions give separately
\begin{align}
 K-k_L
 &=2-w_1-w_2-R_1-R_2-(1-w_L-R_L)\notag\\
 &=1+2\ell-\Delta_R=2\ell+2.
 \label{s7c:noninterlacing-slot}
\end{align}
Combining the universal extraction \eqref{s7c:universal-extraction}
with the same two-component row of Lemma~\ref{lem:homflyrows}, its two requested degrees
are now $k_L+2\ell-2=K-4$ and $k_L+2\ell=K-2$.
Therefore
\begin{equation}\label{s7c:noninterlacing-extraction}
 \Omega_H-\Omega_L=[a^{K-4}](f_1f_2)-[a^{K-2}](f_1f_2).
\end{equation}
If either Laurent row $f_i$ is zero, both reads are zero immediately.
Otherwise apply Theorem~\ref{thm:floor} to the two one-dissent half
contact carriers, subpolygons of a decomposition of the generic side
(Lemma~\ref{lem:carriers}(i); the hypothesis discharge below); it puts
their nonzero product row at degree at least $K$.
Both displayed coefficients are below this floor.  Thus the returned
newborn-free row is zero.  The auxiliary half selectors may themselves
be zero: they are not multipliers of this full term, and the floor
requires their one-dissent turn patterns, not uniform selectors.



For eligible $T$ in the noninterlacing branch, $T\cup\{x\}$ is a
support since $x$ is adjacent to neither $T$ nor $y$.  Every old
neighbour of $y$ is also a neighbour of selected $x$, and hence is
dominated.  Thus $\{y\}$ is an isolated block; the
statement with $x,y$ exchanged is identical.  If a one-newborn
selector is zero its term vanishes.  Otherwise its owner is uniform,
and Lemma~\ref{cb:singleton} makes the required coefficient zero,
including the zero-polynomial case.  Every one-newborn term is zero.

\par\noindent\emph{The different-block alternative.}\par
If $x,y$ lie in different blocks of the newborn-free
high row, both components must be singletons.  Otherwise an undominated
old common neighbour gives a path of length two joining them; the
common-neighbour property follows directly from
\eqref{s7c:bigon-words-eq}.  Each singleton has polynomial $1$ and
writhe $1$.  Every other high block is old and is matched to
its low block by erasing the finger.  For $S=T\cup\{x,y\}$,
\begin{equation}\label{s7c:different-residual}
 U(S)=U(T)\setminus\{x,y\}.
\end{equation}
The triangle and half-factorization construction in
the two-newborn calculation above now distributes every old contact block
once, with its actual owner, between $L_1,L_2$.  It follows that
\begin{equation}\label{s7c:different-polynomials}
 f_H=f_L=f_1f_2,\qquad
 w_H=w_1+w_2+2,\qquad w_L=w_1+w_2.
\end{equation}
If the common full selector or a spectator multiplier is zero, both
newborn-free terms vanish before the row-tied absolute rotations are
used.  Otherwise the same noninterlacing turn computation gives
\eqref{s7c:noninterlacing-rotation}.  Thus
\begin{align}
 k_L
 &=1-(w_1+w_2)-R_L\notag\\
 &=K-1+(R_1+R_2-R_L)=K-2,\label{s7c:different-low-slot}\\
 k_H&=k_L-2=K-4.\label{s7c:different-high-slot}
\end{align}
If $f_1=0$ or $f_2=0$, \eqref{s7c:different-polynomials} makes both
full rows zero without a degree argument.  Otherwise the actual
one-dissent half floors give product degree at least $K$.  Both slots
$K-2,K-4$ are below it, so both reads vanish separately.  The
one-newborn rows already vanished by Lemma~\ref{cb:singleton}.
The same-block and different-block alternatives exhaust the residual
components containing the newborns.  Therefore
\begin{equation}\label{s7c:noninterlacing-return}
 R_{\rm ret}=0\qquad(\epsilon=0).
\end{equation}

\par\medskip\noindent\emph{Complete hypothesis discharge and assembly.}\par

Every floor polynomial above is that of an actual positive carrier
one-circle diagram supplied by Lemma~\ref{lem:carriers} and
Lemma~\ref{cb:products}, and every carrier to which
Theorem~\ref{thm:floor} is applied is a subpolygon of a decomposition of
a generic polygon (Lemma~\ref{lem:carriers}(i)), the theorem's domain.
For the two half contact carriers that polygon is the generic half
(Lemma~\ref{lem:children}(ii)), at either value of $\epsilon$: the
traversal circle of $\lambda_1$, respectively $\lambda_2$, is the interval
$A$, respectively $B$, closed at the cut $\mathbf m=\mu_M$
(Definition~\ref{def:deletion-halves}: $\lambda_1$ follows
$E_M,\ldots,E_{a-1}$ and closes with $[\mu_a,\mu_M]\subset E_a$;
$\lambda_2$ opens with $[\mu_M,\mu_b]\subset E_a$, follows
$E_b,\ldots,E_{M-2}$ and closes with $E_{M-1}$); $T_i$ is a support of
$\lambda_i$ all of whose selected visits are internal to that interval
(the bijection \eqref{s7c:eligible-bijection}); the oriented
reconnections of Definition~\ref{def:smoothing} at those visits act
inside the interval and leave the cut untouched, so closing the interval
at the cut closes the boundary successor of $T$ into the carrier of the
decomposition $T_i$ of $\lambda_i$ through the corner $\mathbf m$ ---
which is how $L_i$ was formed --- with the internal cycles as the other
carriers of $T_i$ (Lemma~\ref{lem:carriers}(i)); and the self-crossings
of $L_i$ are the owner-restricted set $O_i$ of
\eqref{s7c:owner-restriction} (Lemma~\ref{lem:carriers}(iii)).  When
$\epsilon=1$ the enlarged supports $T\cup\{x\}$, $T\cup\{y\}$ and
$T\cup\{x,y\}$ of the generic side do not take the place of this
discharge: $T\cup\{x,y\}$ is not a support, the newborn chords
interlacing, and the daughter of $T\cup\{x\}$ or of $T\cup\{y\}$ that
carries the turn at $\mathbf m$ also carries the opposite smoothing turn
of \eqref{s7c:one-newborn-turns} and is mixed, while a live half contact
carrier is uniform.  In the noninterlacing case the two half contact
carriers are the carriers of the newborn-free support $T$ enlarged by
the smoothed newborns --- of $T\cup\{x,y\}$ when
$\epsilon=0$, when the third carrier is the crossing-free local contact
triangle and the deletion of the curl $y$ from component~1 of $D_A$ is that
smoothing with the triangle discarded --- which is what the closure at the
cut and the successor words of \eqref{s7c:component-data} record; the
daughters of a singleton row are the carriers of the support enlarged by
the singleton (Lemma~\ref{cb:singleton}); and a crossing-free carrier is a
carrier of its own row's decomposition.  The construction in Theorem~\ref{thm:floor} supplies
the needed clean rounding and curl discs; no separate arc-contraction
or arbitrary-front lift is presumed.  Its turn hypotheses are exhausted
as follows: live interlacing rows have two uniform half contact carriers;
live noninterlacing rows have two one-dissent half contact carriers;
singleton rows have a uniform and a one-dissent daughter; and a
crossing-free carrier uses its separately proved $P=1,w=0$ branch.

A zero selector or spectator which actually multiplies the current full
or target term stops that row before a row-dependent absolute-rotation
or degree identity.  A zero Laurent row is checked before every
minimum-degree or coefficient-floor argument.  Independently valid
geometric rotation identities do not require nonzero polynomials,
and coefficient extraction itself is always total.  In particular an
auxiliary one-dissent half or daughter can have zero canonical selector:
it is not a multiplier in the current row, and its actual turn pattern
is precisely a permitted floor hypothesis.  No spectator, selector,
polynomial or coefficient has been divided out anywhere in the proof.

For a bigon, \eqref{s7c:sector-split},
\eqref{s7c:interlacing-return} and \eqref{s7c:noninterlacing-return}
give, with all factors already bound,
\begin{equation}\label{s7c:bigon-total}
 B+R_{\rm ret}=(1-\epsilon)J+\epsilon J
             =sC(\lambda_1)C(\lambda_2).
\end{equation}
The sliding calculation above gives the same result for sliding.
The nonzero neighbour-side signs are either equal or opposite; these
alternatives are exclusive and exhaustive.  Hence
the vertex--edge law holds on its full stated domain.
The three positions of every old label---internal to $A$, internal to
$B$, or mixed---retain all exterior crossing and threading data.
All support bijections have the explicit inverse union operation, and
every actual spectator occurs once.  There is no restriction to an
unthreaded drawing or to a preferred choice of the initial side.

\end{proof}




\begin{theorem}[empty-cusp zero]\label{thm:C-S5}
At a simple empty cusp, $C(P_{\rm no})=0$.
\status{proved (refereed: bench A, 2026-09-05T15:51:04Z; transcribed from CV thm:zeroanchor(S5) (P-25-2, C006))}
\end{theorem}
\begin{proof}
Let the cusp be at vertex $j$ of an $n$-gon, $n\geq4$, with tail-labelled
edges $E_i=[\mu_i,\mu_{i+1}]$. Choose a generic polygon $P$ sufficiently
near the wall on its no-loop side. In cusp case A, where
$\mu_{j+1}(0)$ lies strictly between $\mu_{j-1}(0)$ and $\mu_j(0)$,
put $(c_1,c_2)=(j,j+1)$. In case B, where $\mu_{j-1}(0)$ lies strictly
between $\mu_j(0)$ and $\mu_{j+1}(0)$, put $(c_1,c_2)=(j-1,j)$.
In both cases the corners $c_1,c_2$ are consecutive, with intervening
edge $E_{c_1}$.

By Lemma~\ref{lem:cusp-sides}(ii), their turns on $P$ satisfy
\begin{equation}\label{s5pr:opposed-turns}
\tau_{c_1}(P)=-\tau_{c_2}(P),\qquad
\tau_{c_1}(P),\tau_{c_2}(P)\in\{-1,1\}.
\end{equation}
The cusp is empty, meaning that the two newborn crossing visits on the
loop side are cyclically adjacent. Lemma~\ref{lem:cusp-sides}(iii)
therefore gives that $E_{c_1}$ has no crossing on $P$. This conclusion
uses the lemma's specified-short-arc argument; adjacency alone is not
silently assigned to one of the two complementary traversal arcs.

Fix any decomposition $S\in\operatorname{Ind}(G_P)$.
In the traversal circle, take the closed directed arc from the visit
of $\mu_{c_1}$ to the next vertex visit $\mu_{c_2}$, tracing the single
edge $E_{c_1}$. Its interior has no crossing visit, by the preceding
paragraph. Its endpoints are not crossing visits either: a generic
polygon has distinct vertices and no vertex on a nonincident edge.
Consequently this arc contains no visit of a selected crossing of $S$.

Oriented smoothing changes the traversal only at selected crossing
visits (Definition~\ref{def:smoothing}). Every directed segment of the
chosen arc, including its passages into the two vertex visits, is
therefore retained with the same direction. More explicitly, subdivide
the traversal circle at all polygon-vertex visits and all selected
crossing visits. At a selected pair smoothing interchanges the two
outgoing continuations; it leaves every other local continuation
unchanged. The arc joining $c_1$ to $c_2$ has no selected subdivision
point, so it remains a directed path in the reconnected graph. Each
vertex of that graph has one incoming and one outgoing edge, hence the
graph is a disjoint union of directed cycles. The path lies on one
such cycle. Thus both original vertices belong to the same subpolygon,
say $Q$, of the decomposition $S$.

At each of these original vertices, smoothing has changed neither of
the two incident edge directions. Both remain actual corners of $Q$,
since their original turn signs in~\eqref{s5pr:opposed-turns} are nonzero.
Their turns on $Q$ are the same as on $P$, so one is left and the other
right. Therefore $Q$ is mixed, and $S$ is not uniform
(Definition~\ref{def:uniform}). This argument also covers
$S=\varnothing$, when the one subpolygon is $P$ itself.

Since $S$ was arbitrary, no decomposition of $P$ is uniform. The sum
defining $C$ in Definition~\ref{def:C} has an empty index set, so
\begin{equation}\label{s5pr:empty-sum}
C(P)=(-1)^{\ell(P)}
\sum_{\substack{S\in\operatorname{Ind}(G_P)\\S\text{ uniform}}}
(-1)^{|S|}\prod_{Q'\text{ subpolygon of }S}c(Q')=0.
\end{equation}
No corner coefficient has been evaluated or divided out.

Finally, opposite turns at these consecutive vertices and absence of
crossings on their connecting edge depend only on the turn word and
the crossing set. Those data are constant on each labelled chamber
(Proposition~\ref{prop:chambers}). Hence the same argument applies to
every polygon in that labelled no-loop chamber. The existence of the
pair and the argument are preserved by cyclic relabelling, so they
also apply to its polygon chamber. This proves $C(P_{\rm no})=0$
without needing a separate polynomial-invariance argument.
\end{proof}

\begin{theorem}[soft theorem for $C$]\label{thm:C-soft}
Let $P$ be generic, $j$ a vertex, $q$ admissible, and $P_\varepsilon$ the
soft insertion of Definition~\ref{def:soft}, with attachment signs
$\chi_\pm$. Then for all sufficiently small $\varepsilon>0$
\[
C(P_\varepsilon)=\frac{\chi_-+\chi_+}{2}\,C(P).
\]
\status{new (round~4: the proof consumes the exported geometry of Lemma~\ref{lem:soft-generic}(i)--(iv), chamber constancy, the carrier correspondence of Lemma~\ref{lem:carriers} and the record bridge (Lemma~\ref{rp:record-polynomial}, Theorem~\ref{lp:core}); adopted from the Codex proposal sm4\_round4\_exports\_v1, P-25-24, F120-C-PROOF, re-read by the author; assembly: Lemma~\ref{lem:soft-generic}(i)--(iv) new (round~4; cleared by bench~A on SM5, 2026-09-06T00:32:58Z), Proposition~\ref{prop:C-chamber} held, Lemma~\ref{lem:carriers} (i)--(iii) refereed and (iv) new (cleared by bench~A on SM4), Lemma~\ref{lem:corner-values} held, Lemma~\ref{rp:record-polynomial} refereed, Theorem~\ref{lp:core} new (its identification paragraph refereed on SM4), Lemma~\ref{lem:rot} proved; F-25-120; round~5: ``Gauss word'' for ``crossing word'' in the proof, F-25-148)}
\end{theorem}
\begin{proof}
Write $M=\mu_j$, $M_\varepsilon=M+\varepsilon q$,
$u=\lt_{j-1}$, $v=\lt_j$, $D_\varepsilon=v-\varepsilon q$
and $\tau=\tau_j(P)$. Lemma~\ref{lem:soft-generic} supplies
one initial generic interval. By Proposition~\ref{prop:C-chamber},
$C$ is constant on that chamber, so it suffices to prove the identity
on a smaller initial interval. All choices below can hold
simultaneously because there are finitely many supports and carriers.

\emph{Carrier records and coefficients.}
Use the marked successor model and inherited orientations of
Lemma~\ref{lem:carriers}. For corresponding carriers with the same
retained crossing visits, the cyclic successors and pairings agree.
Each carrier crossing germ follows its containing edge of its own
ambient polygon in the inherited direction, up to a positive scalar
multiple.
Lemma~\ref{lem:soft-generic}(iii) preserves the determinant sign
of each ordered pair of those directions. By
Definition~\ref{def:positive-lift}, the same corresponding visit
is therefore over, and every crossing sign is positive. These
are isomorphic named records, including the correspondence of
oriented crossing-free carriers when the visit sets are empty.
Lemma~\ref{rp:record-polynomial} equates their source polynomials;
the common Laurent substitution and Theorem~\ref{lp:core} identify
these with equal $H^+$-polynomials. If the corresponding rotations
also agree, their crossing counts $m_Q$, rotations $r_Q$ and
$d_Q=1-m_Q-|r_Q|$ agree. Definition~\ref{def:C} then gives equal
corner coefficients $c(Q)$. We now establish each needed carrier
correspondence and rotation assertion.

\emph{Same-sign sector: $\chi_-=\chi_+=-\tau$.}
The inherited Gauss words agree by Lemma~\ref{lem:soft-generic}(iv),
so the independent supports are precisely the same sets $S$.
In the successor cycles of Lemma~\ref{lem:carriers}(i), contracting
the inserted ordinary vertex $M_\varepsilon$ identifies each
carrier with its parent carrier. Exactly the carrier through $M$
gains that vertex. Its old turn $\tau$ is replaced by two turns
$\tau,\tau$, by Lemma~\ref{lem:soft-generic}(ii). All other
original and smoothing corner signs persist by clauses (ii),(iii)
and Lemma~\ref{lem:carriers}(ii). Thus uniformity is equivalent
on the two sides of this correspondence.

For rotation, let $\theta$ be the old principal turn from $u$ to
$v$, and let $\alpha$ be that from $u$ to $q$. The same-sign
determinant conditions put $q$ in the oriented turn wedge:
$0<\alpha<\theta<\pi$ if $\tau=1$, and
$-\pi<\theta<\alpha<0$ if $\tau=-1$.
Indeed the two strict half-plane inequalities choose precisely
the intersection between the rays of $u,v$ inside their principal
angular sector. The two new principal turns tend to $\alpha$
and $\theta-\alpha$, whose sum is $\theta$, with no wrap.
Every other carrier corner has directions converging to its old
regular directions, including smoothing corners involving the
return edge and the successor corner at $B=\mu_{j+1}$.
The full turn sum therefore converges to its parent turn sum.
Lemma~\ref{lem:carriers}(ii) gives regularity of the nearby
carriers, and Lemma~\ref{lem:rot}(i) makes their rotations
integers. Convergence to the old integer implies equality for
all sufficiently small parameters. This applies to every carrier,
including geometrically moving carriers that do not pass through $M$.

Their retained visit records agree by the successor correspondence
and Lemma~\ref{lem:soft-generic}(iii). The preceding coefficient
comparison gives equal products of $c(Q)$. Support sizes do not
change, while $\ell(P_\varepsilon)=\ell(P)+[\tau=1]$.
Substituting in Definition~\ref{def:C} gives
$C(P_\varepsilon)=(-1)^{[\tau=1]}C(P)=-\tau C(P)$,
the required multiplier in this sector.

\emph{Mixed sector: $\chi_-\ne\chi_+$.}
The soft edge has no crossing by Lemma~\ref{lem:soft-generic}(i).
Every selected reconnection therefore leaves the directed path
from $M$ to $M_\varepsilon$ intact. These are consecutive corners
of one carrier, and their turns $-\chi_-$ and $-\chi_+$ are
opposite. No support is uniform, so Definition~\ref{def:C} gives
$C(P_\varepsilon)=0$, the required mixed-sector value.

\emph{Loop sector: $\chi_-=\chi_+=\tau$.}
Let $y$ be the newborn crossing. Its visits are adjacent in the
Gauss word by Lemma~\ref{lem:soft-generic}(iv), so it interlaces
no crossing. The independent supports are exactly $S$ and
$S\cup\{y\}$, with $S\in\Ind(G_P)$.
For a support omitting $y$, Lemma~\ref{lem:carriers}(iii) puts
both visits of $y$ on one carrier $Q'$. Its short arc through
$M,M_\varepsilon$ has no other crossing visit; hence $y$ remains
isolated in the interlacement graph of $Q'$. If the decomposition
is not uniform its term is excluded. If it is uniform,
Lemma~\ref{lem:corner-values}(ii) gives $c(Q')=0$.
Thus every support omitting $y$ contributes zero, including those
whose corresponding parent support was not uniform.

For $S\cup\{y\}$, process $y$ first, using order independence
in Lemma~\ref{lem:carriers}(i). Let $a$ be its visit on the incoming
edge and $b$ its return-edge visit. The clean short arc is
$a\to M\to M_\varepsilon\to b$. Swapping the two outgoing
successors produces the triangle cycle $(b,M,M_\varepsilon)$
and the residual cycle containing $a$. The triangle smoothing
turn follows $D_\varepsilon$ into $u$ and has sign $-\tau$;
the turns at $M,M_\varepsilon$ also have sign $-\tau$.
The three noncollinear segments form an embedded triangle,
so Lemma~\ref{lem:corner-values}(i) gives its coefficient $1$.

All old crossing visits lie on the residual cycle. It replaces
the old corner $M$ by the smoothing corner $y$, whose incoming
and outgoing directions are $u,D_\varepsilon$, of turn sign
$\tau$ by Lemma~\ref{lem:soft-generic}(ii).
The remaining swaps $S$ act only on that residual cycle and give
carriers in bijection with the carriers of $S$ on $P$.
Their original and smoothing corner signs agree; at the distinguished
carrier, the corner $M$ is replaced by $y$ of the same sign.
Consequently $S\cup\{y\}$ is uniform exactly when $S$ is uniform.
No (G1) claim about the intermediate residual polygon is needed:
all resulting curves are carriers of the generic $P_\varepsilon$.

As $y\to M$ and $D_\varepsilon\to v$, every residual principal
turn converges to its old regular counterpart. Integer-valuedness
from Lemma~\ref{lem:rot}(i) again gives equality of all rotations
on a common small interval. Their retained crossing successors and pairings agree under the
inherited correspondence, and the corresponding directed germs
retain their ordered determinant signs,
so the earlier named-record argument gives equal positive-lift
polynomials and corner coefficients. This also covers all other
corresponding carriers, whose positions may have moved.

Every surviving term thus gains exactly one selected crossing and
one triangle coefficient $1$. The original left-turn count changes by
$\ell(P_\varepsilon)-\ell(P)=-[\tau=1]+2[\tau=-1]$.
Definition~\ref{def:C} gives
\[
C(P_\varepsilon)
=(-1)^{-[\tau=1]+2[\tau=-1]+1}C(P).
\]
For $\tau=1$ the exponent is zero; for $\tau=-1$ it is three.
The multiplier is therefore $\tau$, as required. The three sector
multipliers are exactly $(\chi_-+\chi_+)/2$, completing the proof
on the unchanged domain.
\end{proof}

\begin{proposition}[reversal]\label{prop:C-reversal}
$C(\overline P)=(-1)^nC(P)$. \status{new (round~3: the reversal invariance of $H^+_Q$ taken from Theorem~\ref{cf:thm-carrierfloor}(R) and Theorem~\ref{lp:core}, F-25-54; Codex U-P4, P-25-12, adapted)}
\end{proposition}
\begin{proof}
Reversal preserves the crossing set, interlacement, decompositions,
subpolygons and their crossings; it reverses every strand direction, so the
positive lift is preserved (both $u_{\rm o},u_{\rm u}$ change sign and
$\det$ is unchanged), $H^+_Q$ is unchanged, since the positive lift of the
reversed carrier is the reversed diagram of the positive lift of $Q$,
Theorem~\ref{cf:thm-carrierfloor}(R) gives $P_{-D}=P_D$ for every oriented
knot diagram $D$, and Theorem~\ref{lp:core} identifies $P$ with $H$; $m_Q$
and $|r_Q|$ are unchanged,
uniformity is unchanged, and $\ell(\overline P)=n-\ell(P)$
(Lemma~\ref{lem:shift}(iii)). The prefactor gives the sign $(-1)^n$.
\end{proof}

\subsection{The triple wall}

\begin{hypothesis}[R]\label{hyp:R}
At every simple triple wall, $C(P_+)=C(P_-)$. \status{hyp}
\end{hypothesis}

\begin{lemma}[the two triple types]\label{lem:triple-types}
At a simple triple wall at $\{e,f,g\}$ with edge vectors $u_e,u_f,u_g$, put
$\delta_1=\sgn\det(u_e,u_f)$, $\delta_2=\sgn\det(u_e,u_g)$,
$\delta_3=\sgn\det(u_f,u_g)$. Exactly one of the following holds.
\begin{enumerate}
\item[(i)] \emph{Transitive}: $(\delta_1,\delta_2,\delta_3)\notin\{(+,-,+),(-,+,-)\}$.
Then the relation ``$E_i$ passes over $E_j$'' of the positive lift is a total
order on the three strands, and for every subpolygon $Q$ retaining all three
crossings the positive lifts on the two sides differ by a Reidemeister III
move.
\item[(ii)] \emph{Cyclic}: $(\delta_1,\delta_2,\delta_3)\in\{(+,-,+),(-,+,-)\}$.
Then each strand passes over the next cyclically. For a subpolygon
retaining all three crossings, the two sides have that same cyclic
over/under pattern and opposite local triangle configurations.
\end{enumerate}
\status{new (round~4: the subpolygon clauses proved --- common event disc, carrier correspondence across the wall by suppression and reinsertion of the six local visits, exterior planar isotopy of the one carrier, local tournament; adopted from the Codex proposal sm4\_round4\_triple\_star\_v1 (seal 9a82b59a\ldots), F122-PROOF, re-read by the author; F-25-122; round~5: the object the exterior isotopy perturbs from named as the carrier of $S$ at the nongeneric centre, F-25-122)}
\end{lemma}
\begin{proof}
Write $T=\{x_{ef},x_{eg},x_{fg}\}$ for the three crossing labels.
At the central triple point $q$, the three edge directions are
pairwise nonparallel: otherwise two supporting lines coincide,
forcing three original vertices to be collinear, contrary to
$Z_{\rm pt}=\varnothing$.
The three direction determinants are therefore nonzero and retain
their signs near the centre.

Choose a closed event disc about $q$ whose intersection with the
central polygon consists only of these three strand segments, with
their endpoints on its boundary. The point $q$ is interior to the
three edges and is not an original vertex. A fourth edge through
$q$ would be remote from the first three, since adjacent edges under
(G1) meet only at their common endpoint; it would give another
triple in $Z_{\rm c}$, a contradiction. The other finitely many
compact edges thus have positive distance from $q$.
Choose the disc radius small enough to exclude them and all original
vertices, with its boundary transverse to the three strands.
After shrinking the parameter interval, the disc still contains
exactly these three strand segments, all three crossings of $T$,
and no original vertex or other crossing.
Lemma~\ref{lem:triple-sides} gives adjacent pairs of local visits.
Lemma~\ref{lem:wall-sides}, using the sign-change condition in
Definition~\ref{def:walls}(T), reverses each local pair across
the wall and preserves every other visit order.

Let $Q_-$ be any subpolygon on one side retaining all three
crossings, and fix its decomposition $S$.
Definitions~\ref{def:decomposition} and~\ref{def:smoothing}
give $S\cap T=\varnothing$.
Delete the six visits of $T$ from both parent traversal lists.
The remaining cyclic lists of all vertices and other crossing visits
are identical. In particular all selected pairs have the same cyclic
order, so $S$ is independent on the other side as well.

Apply the successor description of Lemma~\ref{lem:carriers}(i).
Every deleted visit is unselected, hence has its ordinary outgoing
traversal arrow. Suppressing these visits commutes with all selected
successor swaps: following the unchanged arrows until the next
remaining mark has the same endpoint on both sides.
The resulting permutations on the remaining marks are identical,
and their cycles identify the carriers on the two sides.
No entire carrier is suppressed: a cycle with only unselected
visits follows the original traversal and reaches an original vertex,
which has not been deleted.
Reinserting each local pair merely subdivides the same oriented arc
of the corresponding cycle, with that pair in the reversed order.
The corresponding carrier $Q_+$ therefore retains all three
crossings. The disc has no selected smoothing site or original
vertex, so both carriers contain all three entire local strand
segments with their inherited orientations.

Lemma~\ref{lem:carriers}(ii),(iii) makes these valid carrier
diagrams on both punctured sides. Each local directed germ is a
positive multiple of its parent edge direction. By
Definition~\ref{def:positive-lift}, its strand from $E_i$ passes
over that from $E_j$ exactly when $\det(u_i,u_j)>0$.
Thus either carrier's local tournament is the parent's determinant
tournament.

For this one carrier, the exterior diagrams may be identified by
planar isotopy. Regard its retained crossings as four-valent vertices,
its original and smoothing corners as two-valent vertices, and the
six transverse event-boundary intersections as boundary vertices.
At a selected crossing this carrier has only one corner, by
Lemma~\ref{lem:carriers}(i); corners belonging to other carriers
are not part of this exterior graph. Away from the event disc, the
graph of the central carrier --- the carrier of $S$ traced on the central
polygon $P(0)$, of which the carriers on both sides are small perturbations
--- has distinct such vertices, noncollapsed arcs and no
degenerate incidence: the only parent triple is inside the disc,
(G1) excludes vertex incidences, and all other intersections are transverse, with distinct crossing
parameters on each edge. Choose disjoint small vertex
neighbourhoods, thin strips around the intervening embedded arcs,
and collars of the six boundary endpoints. Their cyclic attachment
orders persist under a sufficiently small perturbation.
Straightening in these neighbourhoods, strips and collars gives an
ambient planar isotopy of the disc exterior, preserving its boundary
setwise while moving the six marked endpoints to their counterparts.
The nonzero direction determinants at persistent crossings preserve
the physical positive over/under choices during this identification.
Extending the boundary identification across the disc only identifies
coordinates; the two interior triangle drawings are still compared
by the local move below.

A tournament on three strands is a total order unless it is a
directed cycle. It is cyclic exactly when
$\det(u_e,u_f)$, $\det(u_f,u_g)$ and $\det(u_g,u_e)$ have one
sign, equivalently when $(\delta_1,\delta_2,\delta_3)$ is
$(+,-,+)$ or $(-,+,-)$.
In all other cases call the three strands top, middle and bottom
according to their total order. The local triangle interchange has
the top strand over both others and the bottom strand under both
others, so it is the oriented Reidemeister III move.
This proves (i) for the corresponding carrier diagrams.

In a cyclic case the determinant signs still persist, so the cyclic
over/under pattern is the same on both sides. The three reversed
local visit pairs give opposite triangle configurations.
This proves (ii), with no cyclic knot-invariance conclusion.
\end{proof}

\begin{remark}[Historical terminology]
The cyclic triangle change is conventionally called a $\Delta$-move
\cite{MurakamiNakanishi}. This terminology and its general knot-type
behavior are not used as premises here. The preceding lemma asserts
only the explicitly described local crossing pattern in its cyclic case;
it supplies no cyclic knot-invariance theorem.
\end{remark}

\begin{proposition}[what is proved about the triple wall]\label{prop:R-partial}
\begin{enumerate}
\item[(i)] For a transitive triple wall and a decomposition $S$ that contains
none of the three crossings and in which all three lie on one subpolygon,
the term of $S$ is the same on both sides.
\item[(ii)] Across the wall the interlacement of every pair among the three
crossings toggles and no other interlacement changes; the local
interlacement graph on the two sides is either $K_3$ and the empty graph, or
a path and its complement (an edge plus an isolated vertex).
\item[(iii)] Under the hypothesis that no crossing outside the three
interlaces any of them (\emph{outside isolation}; fix an outside decomposition
$Q$), the sum of the terms of all decompositions $Q\cup J$, $J$ a subset of
the three crossings, is the same on both sides, for both local graph orbits.
\end{enumerate}
\status{proved (refereed: bench A, 2026-09-06T00:35:51Z, 2026-09-06T04:22:56Z and 2026-09-06T05:12:19Z; transcribed from RA R-LOC-2, R-PAR-v6, R-EXTERIOR-1 and the proof files RA calls its four accompanying core proofs (six files) (C042, C051, C064; F-25-21); round~3: the proof wording at the floor application corrected, F-25-98 (Codex G-P1, P-25-15); round~4: one citation in the proof narrowed from Literature input~\ref{lit:homfly} to Theorem~\ref{lp:core} (seat L-1/L-5, C4); round~5: the two appeals to Theorem~\ref{thm:floor} reworded to its domain (subpolygons of a decomposition), attributing no geometric hypothesis to Lemma~\ref{lem:carriers}, F-25-132; the disclaimer of Hypothesis~R no longer cites its label, so that the ordering instrument draws no consumption edge from it, F-25-160)}
\end{proposition}
\begin{proof}
The proof transcribes RA's establishing warrants and its six
files whose names end in \texttt{PROOF.md}. The warrants themselves
call their collection ``four accompanying core proofs''; that source
designation and the present file count are distinct.

We use the actual carrier sum of Lemma~\ref{lem:C-X1}. For a carrier
$L$ write $P_L=H_L^+$, $w_L=m_L$, $R_L=|\rot(L)|$,
$d_L=1-w_L-R_L$, and $\omega_L=[a^{d_L}z^0]P_L$.
The term of a support is the product of $\operatorname{wt}(L)\omega_L$
over its carriers. All diagrams below are actual carrier diagrams or
their explicitly specified switched or smoothed diagrams. Products of
owned block polynomials are justified by Lemma~\ref{cb:products};
equal decorated records give equal polynomials by
Lemma~\ref{rp:record-polynomial}. This does not identify the present
definition with an unproved frozen state sum.

\par\medskip\noindent\emph{Local geometry and (ii).}\par
Lemma~\ref{lem:triple-sides} gives a fixed crossing set and three
adjacent pairs of local visits. Lemma~\ref{lem:wall-sides} gives their
reversal across this simple wall and fixes all other visit orders.
Choose the event disc small enough to contain only the three strand
segments: its central point is interior to them, no vertex lies there,
and another strand through that point would violate simplicity.
The three shrinking segments between local visits contain neither a
vertex nor another crossing.

An adjacent transposition of visits of two distinct crossings changes
only their mutual interlacement. Indeed it moves one visit past only
the other; every other pair keeps its four-point cyclic order. For the
two named crossings, exactly one endpoint passes through the boundary
of the interval cut by the other chord, so alternation toggles.
There are three such transpositions, one for each local pair. Thus the
local graph is complemented and every other edge is unchanged. On
three vertices the complementary graphs have edge counts $0,3$ or
$1,2$; the latter graphs are an edge and a path. This proves (ii).

\par\medskip\noindent\emph{Rotations and (i).}\par
Every selected crossing that persists gives two corners with nonzero
direction determinants near the centre. Polygon corners likewise
remain nonzero by (G1). Carriers with fixed successor patterns and no
collapsed edge between selected local crossings are therefore regular
families through the wall. Their individual directions may change.
Their principal turns vary continuously without meeting $\pm\pi$;
their integer rotation numbers are constant by Lemma~\ref{lem:rot}.
Their corner signs and weights are constant as well.
For the support in (i), no local crossing is selected. The carrier
containing the triple undergoes the transitive oriented RIII change
of Lemma~\ref{lem:triple-types}(i). Theorem~\ref{lp:core} --- Reidemeister-III
invariance of the local polynomial, from Literature input~\ref{lp:lm}, and its
identification with $H$ --- therefore preserves its polynomial. It has the same retained crossing
labels, hence the same writhe; the preceding regular-family argument
preserves its rotation and weight. Every other carrier has the same
decorated record, writhe, rotation and weight. Multiplying their
coefficient reads proves (i), also when a factor is zero.

\par\medskip\noindent\emph{Availability and the outside factor.}\par
For (iii) fix the stated outside independent support $Q$ and write
$T$ for the three local crossings. We prove the identity already under
the weaker premise that \emph{every selected member of $Q$} interlaces
no member of $T$. The printed premise of (iii), which quantifies over
every outside crossing, implies this premise. We do not change that
printed statement. RA uses the weaker, support-dependent meaning of
full availability.

For any outside crossing $y$, its two visits avoid the interiors of
the three adjacent two-visit clumps. The two arcs cut out by $y$ thus
colour the three whole clumps in two colours. A member of $T$
interlaces $y$ exactly when its two clumps have different colours.
Three objects have either zero or two bichromatic pairs. Consequently
the local mask $I(y)=N(y)\cap T$ has size zero or two. It is fixed
across the wall by (ii). More generally,
\[
 \operatorname{avail}(Q)=T\setminus\bigcup_{y\in Q}I(y)
\]
has size three, one or zero, since the union of distinct two-element
subsets of a three-element set is the whole set. Only size three is
used in the proof that follows.

Smooth $Q$ first. A resulting carrier containing no local visit,
including a local smoothing-site visit, survives any subsequent
independent local smoothings arc for arc. Conversely undoing those
local reconnections on a triangle-disjoint carrier changes no half-edge
of that carrier. This is a bijection of the exterior carriers in all
rows. Their self-crossing sets also stay fixed. To see this, an
undominated exterior crossing adjacent to an available member of $T$
would belong to its connected block and hence to its owner, by
Lemma~\ref{cb:products}. That exterior owner would contain local
visits, a contradiction. Adding local selections removes only those
labels and their neighbours, by~\eqref{cb:greedy-step}; it therefore
removes no exterior self-crossing.

Erasing the six local visits identifies these exterior successor cycles
across the wall. Their decorated records, positive writhes and corner
signs agree; the regular-family argument gives equal rotations.
Their polynomials and coefficient slots agree as well. Thus all rows
have one common exterior factor $C_Q$, which may be zero. For a side
$\nu$ write $T_\nu(J)$ for the \emph{complete} term of $Q\cup J$,
setting it to zero when that support is absent. We only multiply by
$C_Q$; we never divide by it.

Each member of $Q$ has mask zero. Its visits lie in one of the three
gaps between the local clumps: the three gaps have distinct vectors
of sides relative to the three local chords. Smoothing $Q$ therefore
acts separately inside each gap. A printed gap below means its
boundary-to-boundary successor path after those smoothings; internal
closed successor cycles are precisely the exterior carriers just
handled. This supplies the same local word calculation for arbitrary
$Q$, not just for $Q=\varnothing$.

\par\medskip\noindent\emph{Sign classification.}\par
In cyclic strand order put $a=(u_1,u_2)$, $b=(u_1,u_3)$,
$c=(u_2,u_3)$, and let $s_a,s_b,s_c$ be their determinant signs.
Writing $\delta$ for the sign of the three-line determinant, Cramer's
rule for the intersection parameters gives the three order signs
\begin{equation}\label{rtr:line-orders}
 (q_1,q_2,q_3)=-\delta(s_as_b,s_as_c,s_bs_c).
\end{equation}
For example, if $t_{ij}$ is the parameter on line $i$ of its
intersection with line $j$, the numerator of
$t_{12}-t_{13}$ over the common denominator
$\det(u_1,u_2)\det(u_1,u_3)$ is the negative three-line determinant.
The other two rows follow by expanding $t_{21}-t_{23}$ and
$t_{31}-t_{32}$ with the same fixed determinant order.
The local graph indicators, read from the three consecutive clumps,
are respectively $q_1=-1$, $q_2=+1$, $q_3=-1$ for $ab,ac,bc$.
They all agree precisely when $s_a=s_c=-s_b$. These are the two
extreme sign branches; the other six are generic.

A pair sharing a strand $u$ is called separating when the other two
directions lie on opposite sides of $u$. For pairs $ab,ac,bc$ this
requires, respectively, $s_a=-s_b$, $s_a=s_c$, $s_b=-s_c$.
In a generic branch exactly one holds. Equation~\eqref{rtr:line-orders}
identifies that pair as the independent pair complementary to the
path's middle vertex. These statements follow in each of the three
sign pairs
\[
 \begin{array}{c|c|c}
 (s_a,s_b,s_c)&\text{separating pair}&\text{other pairs}\\ \hline
 +++\text{ or }---&ac&ab,bc\\
 ++-\text{ or }--+&bc&ab,ac\\
 +--\text{ or }-++&ab&ac,bc
 \end{array}
\]
and do not presume a symmetry of the outside gaps.

\par\medskip\noindent\emph{Generic orbit: words and common rows.}\par
Choose the ordering of the three directions from their transitive
over-order or its reverse that is even relative to their cyclic
traversal order. A cyclic relabelling then puts the separating pair
at $ac$ and gives $s_a=s_b=s_c=\sigma$, where $\sigma=\pm1$.
No strand orientation is reversed. The words on the path and edge
sides are
\begin{equation}\label{rtr:generic-words}
 P=a b A a c B b c C,\qquad E=b a A c a B c b C.
\end{equation}
The gaps $A,B,C$ are carried along with this relabelling. The entire
support and local-undominated table is
\[
\begin{array}{c|ccccc}
 P:J&\varnothing&a&b&c&ac\\
 U(Q\cup J)\cap T&abc&c&\varnothing&a&\varnothing\\ \hline
 E:J&\varnothing&a&b&c&ab,bc\\
 U(Q\cup J)\cap T&abc&b&ac&b&\varnothing
\end{array}
\]
where $ac$ is connected in the $E,b$ row. All unlisted supports are
absent. The empty row's contact carrier has successor string $ABC$
on both sides. Its local lift has transitive over-order and undergoes
ordinary RIII. Its polynomial, writhe, rotation and weight therefore
agree exactly as in (i); no bijection of the changing local block
partition is required. Hence $T_P(\varnothing)=T_E(\varnothing)$.

For endpoint $a$ the carriers are $BC\mid A$, with residual local
word $cBcC$ on $P$ and $bBbC$ on $E$. Outside survivors after $a$
have mask zero or $bc$. Thus $b,c$ are twins with respect to those
survivors; fixing all outside labels and replacing $c$ by $b$ identifies
the two undominated graphs and their owners. This also identifies
their decorated records: the relocated first visit changes from the
$u_2$ strand to $u_1$, the second remains on $u_3$, and
$\sgn\det(u_2,u_3)=\sgn\det(u_1,u_3)=\sigma$ preserves which visit
is over. Both crossing signs are positive. The two smoothing turns
are $u_1\to u_2$ on $BC$ and $u_2\to u_1$ on $A$.
The carrier families are regular through the wall, so their rotations,
weights, writhes, polynomials and slots agree.
Consequently $T_P(a)=T_E(a)$.

For endpoint $c$ the carriers are $AC\mid B$, and the local residual
words are $aAaC$ and $bAbC$. Outside survivors have mask zero or
$ab$. Replacing $a$ by $b$ fixes their graph incidences and owners.
The first visit stays on $u_1$, while the second changes from $u_2$
to $u_3$; $\sgn\det(u_1,u_2)=\sgn\det(u_1,u_3)$ preserves the
over/under bit. The local turns are $u_2\to u_3$ and $u_3\to u_2$.
The same regular-family and decorated-record argument therefore proves
$T_P(c)=T_E(c)$. This verifies each endpoint separately.

Each of the other two independent pair supports is selector-zero.
For such a pair the shared strand $u$ is nonseparating, so
$\sgn\det(u,v)=\sgn\det(u,w)$ for the other directions. Its short
arc between local visits is intact after smoothing $Q$. Smoothing
the two local crossings puts both ends on one carrier, with turns
$v\to u$ and $u\to w$, or with $v,w$ interchanged. Those turn
signs are opposite. This carrier is mixed whatever happens outside
the disc; hence $T_E(ab)=T_E(bc)=0$.

\par\medskip\noindent\emph{Generic selected singleton and pair.}\par
The remaining successor rows are
\[
 P,b:C\mid AB,\qquad E,b:C\mid AB,\qquad P,ac:C\mid A\mid B.
\]
The $AB$ carrier in $E,b$ has two extra positive crossings, with word
$aAc aBc$ (spaces are immaterial); its $P,b$ mate has neither.
The corner ledger is
\begin{equation}\label{rtr:generic-turns}
\begin{array}{c|cc}
 &C&\text{other carriers}\\ \hline
 P,b\text{ and }E,b&\sigma&AB:-\sigma\\
 P,ac&\sigma,\sigma&A:-\sigma,\ B:-\sigma
\end{array}
\end{equation}
Every nonlocal corner partitions according to the displayed successor
strings. The actual weight definition gives
\[
 \operatorname{wt}(C_{ac})=-\sigma\operatorname{wt}(C_b),\qquad
 \operatorname{wt}(A)\operatorname{wt}(B)
       =\sigma\operatorname{wt}(AB).
\]
For left turns an extra corner changes $(-1)^{\text{corner count}}$;
for right turns the weight remains one. If inherited signs make one
side mixed, they make the matching product mixed as well. Thus the
pair's full selector is the negative of the singleton's, including
zero, and the two singleton selectors are equal.

The principal-angle formula gives
$\rot(C_{ac})=\rot(C_b)$ and
$\rot(A)+\rot(B)=\rot(AB)$. To verify the first equality, the
transitive directions lie in a common open half-plane, and the turn
from the first to the last equals the sum of the two turns through
the middle, all on the same principal branch. The remaining local
turns are the negatives of these, giving the second equality.
The $P,b$ and $E,b$ rotations agree by their regular families.
If the common singleton selector is zero all three terms vanish.
Otherwise $A,B$ are uniform of sign $-\sigma$, so
Lemma~\ref{lem:uniformrot} gives
$R_A+R_B=R_{AB}$.

Let $D_H$ be the actual $E,b$ diagram on $AB$ and $D_L$ its $P,b$
mate. Switching $a$ in $D_H$ makes $a,c$ an opposite-sign oriented
RII pair. Its short boundary arcs, including the selected $b$
connector, lie in the clean event disc. No gap arc or spectator is
in its lens. Deleting this bigon gives $D_L$, with its whole outside
record. Put $D_0=\operatorname{smooth}_aD_H$. The skein equation
and the two additional positive crossings give
\begin{align}
 P_{D_H}&=a^{-2}P_{D_L}+a^{-1}zP_{D_0},\notag\\
 w_H&=w_L+2,\qquad d_H=d_L-2,\notag\\
 \omega_H-\omega_L&=[a^{d_L-1}z^{-1}]P_{D_0}.
 \label{rtr:generic-skein}
\end{align}
The last line is extraction at $a^{d_H}z^0$; its first summand
extracts $a^{d_H+2}=a^{d_L}$.

The two components of $D_0$ have strings $A,B$. In the full $E$
word their side vectors relative to $(b,a,c)$ are
$A=(1,1,0)$, $B=(1,0,1)$, $C=(0,0,0)$.
Outside survivors after $b$ therefore have mask zero or $ac$.
Mask-zero self-crossings have exactly the component records of the
$P,ac$ carriers. A mask-$ac$ crossing is mixed in $D_0$ and absent
from those pair-row records. The local crossing $c$ is mixed as well.
Thus the intrinsic component polynomials are precisely $P_A,P_B$,
also when a component has no self-crossings; this follows from
Lemma~\ref{rp:record-polynomial} with the actual records just given.
Put $f_i=[z^0]P_i$, $D=d_A+d_B$, and let $\ell$ be the linking
number of this actual auxiliary link. All its mixed crossings are
positive, so their number is $2\ell$. Removing the one local mixed
crossing, which is absent from $D_L$, gives
\begin{align}
 w_L&=w_A+w_B+2\ell-1,\notag\\
 D&=2-(w_A+w_B)-(R_A+R_B)\notag\\
  &=2-(w_L-2\ell+1)-R_{AB}\notag\\
  &=d_L+2\ell.\label{rtr:generic-slot}
\end{align}
Lemma~\ref{lem:homflyrows} gives the actual mixed-link row, without
requiring that link to be split. Equations~\eqref{rtr:generic-skein}
and~\eqref{rtr:generic-slot} therefore yield
\begin{align}
 \omega_H-\omega_L
 &= [a^{d_L-2+2\ell}]f_Af_B
       -[a^{d_L+2\ell}]f_Af_B\notag\\
 &= [a^{D-2}]f_Af_B-[a^D]f_Af_B
  =-\omega_A\omega_B.\label{rtr:generic-result}
\end{align}
For the last equality, if either row is zero all these coefficients
are zero. Otherwise the actual uniform carriers are subpolygons of a
decomposition of the generic side (Lemma~\ref{lem:carriers}(i)), the
domain of Theorem~\ref{thm:floor}, so their exponents are bounded below by
$d_A,d_B$. No product exponent is below $D$, and its coefficient at
$D$ is the product of the two coefficients at those bounds.
Crossing-free carriers are included in that theorem and have
polynomial one. The common $C$ carrier has the same intrinsic record,
writhe and rotation in all three rows; its coefficient read is common.
Multiplying~\eqref{rtr:generic-result} by it, the exterior factor and
the singleton selector gives
\begin{equation}\label{rtr:generic-couple}
 T_E(b)=T_P(b)+T_P(ac).
\end{equation}
Together with the common and zero rows, this proves equality of the
two full-availability sums in the generic orbit.

\par\medskip\noindent\emph{Extreme orbit: words and owners.}\par
Relabel the crossings $x=(u_1,u_2)$, $y=(u_1,u_3)$,
$z=(u_2,u_3)$. Here $z$ as a crossing label is distinguished by
context from the HOMFLY variable. The complete and empty local graphs
have words
\begin{equation}\label{rtr:extreme-words}
 H=x y A z x B y z C,\qquad L=y x A x z B z y C.
\end{equation}
The sign classification gives $s_x=s_z=-s_y=\sigma$.
On $L$ every local subset is independent; on $H$ only the empty
set and the three singletons are. The gap side vectors relative to
$(x,y,z)$ are
\begin{equation}\label{rtr:extreme-masks}
 A=(1,1,0),\qquad B=(0,1,1),\qquad C=(0,0,0).
\end{equation}
Thus crossings joining $AB,AC,BC$ have masks $xz,xy,yz$,
respectively. Mask-zero crossings have both visits in one gap.
The legitimate full support $Q\cup\{x,y,z\}$ on $L$ has the three
outer carriers with strings $A,B,C$ and the central triangle. Write
$P_G$, $f_G=[z^0]P_G$, $w_G$ for the actual outer data, $G=A,B,C$.
These symbols always refer to those actual carriers. A crossing-free
one has $P_G=f_G=1$ and $w_G=0$.

The local smoothing turns on the outer carriers are $u_2\to u_1$,
$u_3\to u_2$, $u_1\to u_3$ on $A,B,C$. Each has sign
$s_o=-\sigma$. The opposite three turns make the central triangle
uniform of sign $\sigma$. It is the small triangle inside the event
disc, has three corners and no crossings, and therefore has
$R=1$, slot zero and coefficient one.

\par\medskip\noindent\emph{Extreme pair rows vanish.}\par
In a local pair row on $L$, the remaining local crossing $r$ is
undominated. Any outside neighbour of $r$ has a two-letter mask and
is consequently adjacent to one of the selected pair, so is dominated.
Thus $r$ is an isolated self-crossing on its owner. If that owner is
mixed its weight is zero; if it is uniform
Lemma~\ref{cb:singleton} makes its corner coefficient zero. Each
pair term on $L$ therefore vanishes. The corresponding pairs on $H$
are absent.

\par\medskip\noindent\emph{Extreme singleton transport.}\par
Here are the three successor calculations separately. The $H$ words
retain the visits of dominated local crossings only to identify the
arcs; deleting them gives the common and affected strings.
\begin{equation}\label{rtr:singleton-table}
\begin{array}{c|c|c|c|c|cc}
 j&\text{common}&\text{affected}&H\text{ strings}
   &L\text{ affected word}&q&r\\ \hline
 x&A&BC&yAz\mid ByzC&zBzyCy&y&z\\
 y&C&AB&AzxB\mid zCx&xAxzBz&x&z\\
 z&B&CA&xBy\mid CxyA&yCyxAx&y&x
\end{array}
\end{equation}
The affected carrier has one selected-local corner of sign $\sigma$;
the common carrier has its mate of sign $-\sigma$. This follows
row by row from $u_1\to u_2$, $u_3\to u_1$, $u_2\to u_3$ on
the affected carriers. Their corner lists and weights agree across
the wall. The two carriers in each row form regular families through
it: the unselected local crossings create no corner and the only
selected local turn has nonzero determinant. Their rotations agree.

After deleting the displayed local visits, the common records are
identical. An outside survivor after $j$ has mask zero or the
complementary pair mask, by the parity argument. The latter has one
visit in each affected gap; all these crossings are retained in the
clean affected diagram on $H$ and in its $L$ mate. Call these actual
diagrams $D_0,D_+$, and their common absolute rotation $R$.
The only additional self-crossings of $D_+$ are $q,r$ in the table.
Switching $q$ makes an opposite-sign oriented RII pair. The two
short bigon edges consist of adjacent local strand subarcs and the
already selected $j$ connector; they lie in the clean event disc,
without any gap arc or outside strand in its lens. Deleting the pair
gives exactly $D_0$ with its outside record. Thus, putting
$d_0=1-w_0-R$ and $D_* =\operatorname{smooth}_qD_+$, the skein
equation gives
\begin{align}
 w_+&=w_0+2,\qquad d_+=d_0-2,\notag\\
 P_{D_+}&=a^{-2}P_{D_0}+a^{-1}zP_{D_*},\notag\\
 \omega_+-\omega_0&=[a^{d_0-1}z^{-1}]P_{D_*}.
 \label{rtr:singleton-skein}
\end{align}

The two components of $D_*$, for the three rows of
\eqref{rtr:singleton-table}, are respectively
\begin{equation}\label{rtr:singleton-components}
 C\mid zBz,\qquad A\mid zBz,\qquad C\mid xAx.
\end{equation}
Every complementary-mask retained crossing is mixed; every mask-zero
crossing lies in the indicated gap with exactly its full-support
outer record. The extra self-crossing $r$ on the second component is
isolated: any outside neighbour also meets the selected $q$ and is
dominated in the legitimate support $Q\cup\{j,q\}$.
Lemma~\ref{cb:products} and the single-crossing polynomial one
therefore give component rows $(f_C,f_B)$, $(f_A,f_B)$,
$(f_C,f_A)$ in the three cases. Their component writhes are
$(w_C,w_B+1)$, $(w_A,w_B+1)$, $(w_C,w_A+1)$: the singleton
polynomial is one, but its writhe has not been discarded.

If the matched singleton selector is zero, both complete terms are
zero. Otherwise the affected parent is uniform of sign $\sigma$.
The two additional smoothings $q,r$ split it into the two clean outer
carriers in~\eqref{rtr:singleton-components} and the central triangle.
For $j=x$, $q=y$ puts the $-\sigma$ corner on $C$ and $r=z$ puts
one on $B$. For $j=y$, $q=x$ puts one on $A$ and $r=z$ puts one
on $B$. For $j=z$, $q=y$ puts one on $C$ and $r=x$ puts one on
$A$. Each clean outer carrier has exactly one dissent and at least
two inherited corners, since Lemma~\ref{lem:carriers} gives at least
three corners on every carrier. Each successive smoothing introduces
two opposite principal turns and leaves the other turns unchanged.
Signed rotation is additive by Lemma~\ref{lem:rot}. The one-dissent
rotation result, with reversal if necessary, puts both outer rotations
on the parent's sign ray; the central triangle has that same sign.
Consequently, for the two clean outer absolute rotations,
\begin{equation}\label{rtr:singleton-rotation}
 R_1+R_2+1=R.
\end{equation}
Let $\ell$ be the linking number in $D_*$ and let $w_1,w_2$ be
the \emph{clean full-support} outer writhes, without the extra $r$.
The clean diagram $D_0$ contains the two outer self-crossing sets
and exactly $2\ell$ positive crossings that become mixed in $D_*$.
Thus
\begin{align}
 w_0&=w_1+w_2+2\ell,\notag\\
 D:=d_1+d_2
 &=2-(w_1+w_2)-(R_1+R_2)\notag\\
 &=2-(w_0-2\ell)-(R-1)\notag\\
 &=d_0+2+2\ell.\label{rtr:singleton-slot}
\end{align}
Lemma~\ref{lem:homflyrows} gives
$[z^{-1}]P_{D_*}=a^{-2\ell}(a-a^{-1})f_1f_2$ with the intrinsic
component polynomials just identified. The floor theorem (Theorem~\ref{thm:floor}) applies to
the actual one-dissent outer carriers, subpolygons of a decomposition of
the generic side (Lemma~\ref{lem:carriers}(i)). If a row is zero
the following support assertion is vacuous. Otherwise its exponents
are at least $d_i$, so the product has none below $D$. Extracting
the two terms in~\eqref{rtr:singleton-skein} gives
\begin{align}
 \omega_+-\omega_0
 &= [a^{d_0-2+2\ell}]f_1f_2
       -[a^{d_0+2\ell}]f_1f_2\notag\\
 &= [a^{D-4}]f_1f_2-[a^{D-2}]f_1f_2=0.
 \label{rtr:singleton-result}
\end{align}
Common-carrier and exterior reads agree, as do every weight and
rotation. Multiplication now proves each of
$T_H(x)=T_L(x)$, $T_H(y)=T_L(y)$, $T_H(z)=T_L(z)$.

\par\medskip\noindent\emph{Extreme empty/full couple: two switches.}\par
The empty local support has one contact carrier on each side,
with diagrams $D_H,D_L$, common positive writhe $w$, absolute
rotation $R$, weight $W$, and slot $d=1-w-R$. The common writhe
counts the same retained labels. The unsmoothed local crossings are
not corners, so the regular-family argument supplies the common
rotation and weight. Its row difference is
$C_QW(\omega_H-\omega_L)$.

The original local over-order is cyclic. When $\sigma=+1$ it is
$u_1$ over $u_2$, $u_3$ over $u_1$, $u_2$ over $u_3$; switching
$x$ makes the total order $u_2>u_3>u_1$. When $\sigma=-1$ all
arrows reverse and the switched order is $u_1>u_3>u_2$.
The two switched diagrams are therefore related by ordinary RIII
with the same outside record. Their HOMFLY values agree. Subtract
the two skein equations at $x$; their switched terms cancel and give
\begin{equation}\label{rtr:first-switch}
 P_{D_H}-P_{D_L}
   =a^{-1}z\bigl(P_{D_H^x}-P_{D_L^x}\bigr),
\end{equation}
where the superscript $x$ denotes oriented smoothing, not switching.
After that smoothing the successor rows are
\[
 H:yAz\mid ByzC,\qquad L:A\mid zBzyCy.
\]
Switching $y$ in each leaves $y$ negative and $z$ positive. Their
short local arcs have the orientation-compatible pairing inherited
from the $x$ connector, and bound an empty RII bigon inside the
event disc. Deleting it on both sides leaves the same boundary
pairing $A\mid BC$ and the same outside record. The two switched
link polynomials agree, by RII invariance and
Lemma~\ref{rp:record-polynomial}. Subtracting these two skeins at
$y$ and substituting into~\eqref{rtr:first-switch} yields
\begin{equation}\label{rtr:two-switches}
 P_{D_H}-P_{D_L}
   =a^{-2}z^2\bigl(P_{D_H^{xy}}-P_{D_L^{xy}}\bigr).
\end{equation}
Only matched switched diagrams have been identified: neither is
identified with the opposite \emph{unswitched} diagram.

Smoothing $x$ on either side splits the contact knot. On $H$, $y$
is then mixed, so its smoothing rejoins the two components;
$D_H^{xy}$ is a knot. On $L$, $y$ is a self-crossing and its
smoothing splits again; $J=D_L^{xy}$ has three components
\[
 A,\qquad C,\qquad zBz.
\]
Knot parity in Theorem~\ref{lp:core} gives
$[z^{-2}]P_{D_H^{xy}}=0$. In $J$, the mask-zero crossings give
the full-state outer records, while every two-letter-mask crossing
is mixed. The remaining $z$ is isolated on its component, since an
outside neighbour also meets selected $x$ or $y$. Its polynomial
factor is one and its writhe is one. The intrinsic component rows
are therefore $f_A,f_C,f_B$, by Lemma~\ref{cb:products}.
Write $K=f_Af_Bf_C$ and $\Lambda$ for the sum of the three pairwise
linking numbers of $J$.

For clarity the lowest row needed here follows from the skein
relations themselves. A two-component link has no $z^{-3}$ row:
switch mixed crossings until it is split. Each smoothing is a knot,
whose $z^{-4}$ row is zero; the split endpoint is
$(a-a^{-1})z^{-1}$ times two knot polynomials and also has no
$z^{-3}$ row. The skein relation propagates that zero back.
For three components, switching a positive mixed crossing to a
negative one changes the total linking number by minus one. Its
smoothing has two components and contributes no $z^{-3}$ row.
Thus at row $z^{-2}$ the positive value is $a^{-2}$ times the
negative value; the reverse switch gives the inverse factor.
Switch until the components have a total over-order, and translate
their successive height levels apart to obtain the split endpoint.
The intrinsic component diagrams are unchanged. The split polynomial
is $((a-a^{-1})/z)^2P_AP_BP_C$, by the split product proof of
Lemma~\ref{lem:homflyrows}, applied successively (or
Theorem~\ref{mp:stack} to its three clean one-component blocks).
The net linking change gives
\begin{equation}\label{rtr:three-component-row}
 [z^{-2}]P_J=a^{-2\Lambda}(a-a^{-1})^2K.
\end{equation}
Extracting $a^dz^0$ in~\eqref{rtr:two-switches} now gives
\begin{equation}\label{rtr:extreme-extraction}
 \omega_H-\omega_L
  =-[a^{d+2+2\Lambda}](a-a^{-1})^2K.
\end{equation}

\par\medskip\noindent\emph{Extreme empty/full couple: selectors and floor.}\par
All inherited corners of the empty contact carrier partition among
the three outer carriers. Each outer has its one local corner of
sign $s_o=-\sigma$, while the central triangle has three of the
opposite sign. Let $W_{\rm full}$ be the product of these four
carrier weights. If the parent is mixed, $W=0$. A nonzero full
selector would force all outer carriers to be uniform with sign
$s_o$, hence all inherited parent corners to have that sign, a
contradiction. Both sides of the proposed empty/full identity are
then zero.

Assume the parent uniform and put $\chi=1$ if its sign is $s_o$,
$\chi=0$ if its sign is $-s_o$. If $\chi=1$, all outer carriers
are uniform and $W_{\rm full}=-W$: when $s_o=-1$ the three
outers are all right and the left central triangle contributes
$(-1)^3=-1$; when $s_o=+1$ the outers contain three more left
corners than the parent and the right central weight is one.
If $\chi=0$, each outer has exactly one dissent and at least two
inherited corners, so it is mixed and $W_{\rm full}=0$.

Every smoothing creates two opposite principal turns. Summing all
turns across the three smoothings therefore gives
\[
 \rot(A)+\rot(B)+\rot(C)+\rot(\text{central})
       =\rot(\text{parent}).
\]
In the live branch $\chi=1$ the outer rotations have the parent's
sign and the central triangle has the opposite sign and absolute
rotation one. In the dead branch $\chi=0$ the one-dissent rotation
lemma puts the outer rotations on the parent's sign ray, as does
the central triangle. Hence
\begin{equation}\label{rtr:extreme-rotations}
 R_A+R_B+R_C=
 \begin{cases}R+1,&\chi=1,\\ R-1,&\chi=0.\end{cases}
\end{equation}
The empty parent retains the three positive local crossings, the
outer self-crossings, and all positive mixed crossings of $J$.
There are $2\Lambda$ of the latter. Consequently
\begin{align}
 w&=3+w_A+w_B+w_C+2\Lambda,\notag\\
 D:=d_A+d_B+d_C
   &=3-(w_A+w_B+w_C)-(R_A+R_B+R_C)\notag\\
   &=6-w+2\Lambda-(R_A+R_B+R_C)\notag\\
   &=\begin{cases}d+4+2\Lambda,&\chi=1,\\
                    d+6+2\Lambda,&\chi=0.
      \end{cases}\label{rtr:extreme-slots}
\end{align}
Each actual outer carrier is uniform or exactly one-dissent. Each
is a subpolygon of a decomposition of the generic side
(Lemma~\ref{lem:carriers}(i)), the domain of Theorem~\ref{thm:floor}; with
reversal for negative inherited signs, that theorem bounds the support of
$f_G$ below by $d_G$. If some $f_G=0$
then $K=0$ and all the following coefficients vanish, without any
degree assigned to the zero polynomial. Otherwise no exponent of
$K$ is below $D$ and
$[a^D]K=\omega_A\omega_B\omega_C$.

Expand $(a-a^{-1})^2=a^2-2+a^{-2}$ in
\eqref{rtr:extreme-extraction}. In the live branch its three sampled
degrees of $K$ are $D-4,D-2,D$; only the last can contribute.
In the dead branch they are $D-6,D-4,D-2$, all below the floor.
Thus
\begin{equation}\label{rtr:extreme-result}
 \omega_H-\omega_L=
 \begin{cases}-\omega_A\omega_B\omega_C,&\chi=1,\\
               0,&\chi=0.
 \end{cases}
\end{equation}
The central coefficient is one. In the live branch multiplication by
$C_QW$ and $W_{\rm full}=-W$ gives
\begin{align}
 T_H(\varnothing)-T_L(\varnothing)
 &= -C_QW\omega_A\omega_B\omega_C\notag\\
 &= C_QW_{\rm full}\omega_A\omega_B\omega_C
  =T_L(xyz).\label{rtr:extreme-couple}
\end{align}
The mixed-parent and uniform-dead branches were both zero above.
Together with the three singleton transports and the three pair
zeros, this proves equality of the full-availability sums in the
extreme orbit and completes (iii).

This proof uses neither Hypothesis~R nor the aggregate
vertex--edge conclusion. Each skein triple, component, owner,
rotation and coefficient slot has been specified for this triple
wall. No selector, exterior factor or coefficient has been cancelled.
The availability-one and availability-zero fibres of a general wall
are not proved here; no general triple-wall conclusion follows.
\end{proof}


\begin{remark}\label{rem:R3-maintext}
Hypothesis~R is open in both types: in the transitive type the
decompositions smoothing one or two of the three crossings are reindexed
across the wall (their retained crossings and even their existence change,
by (ii)), and in the cyclic type even the unsmoothed term changes. Exact
computations on hundreds of triple walls, including cyclic ones, have found
the hypothesis satisfied in every case examined; those computations are
evidence, not proof, and are not part of this document's dependency graph.
\end{remark}
